\documentclass[11pt]{article}
\usepackage[margin=1.05in]{geometry}
\usepackage{amsmath,amssymb,amsthm,mathtools}
\usepackage{enumitem}
\usepackage{booktabs}
\usepackage{array}
\usepackage{xcolor}
\usepackage{microtype}
\definecolor{linkblue}{RGB}{20,55,120}
\usepackage[colorlinks=true,linkcolor=linkblue,citecolor=linkblue,urlcolor=linkblue]{hyperref}
\usepackage[capitalize,noabbrev]{cleveref}
\emergencystretch=2em

\newtheorem{theorem}{Theorem}[section]
\newtheorem{lemma}[theorem]{Lemma}
\newtheorem{proposition}[theorem]{Proposition}
\newtheorem{corollary}[theorem]{Corollary}
\newtheorem{fact}[theorem]{Fact}
\theoremstyle{definition}
\newtheorem{definition}[theorem]{Definition}
\newtheorem{conjecture}[theorem]{Conjecture}
\newtheorem{problem}[theorem]{Problem}
\newtheorem{residualinput}[theorem]{Residual Input}
\theoremstyle{remark}
\newtheorem{remark}[theorem]{Remark}

\newcommand{\F}{\mathbb F}
\newcommand{\BB}{\mathbb B}
\newcommand{\RS}{\mathrm{RS}}
\newcommand{\Poly}{\mathrm{poly}}
\newcommand{\eca}{\varepsilon_{\mathrm{ca}}}
\newcommand{\emca}{\varepsilon_{\mathrm{mca}}}
\newcommand{\List}{\operatorname{List}}
\newcommand{\Bad}{\operatorname{Bad}}
\newcommand{\Syn}{\operatorname{Syn}}
\newcommand{\Supp}{\operatorname{Supp}}
\newcommand{\dist}{\Delta}
\newcommand{\distS}{\Delta_S}
\newcommand{\Hb}{H_2}
\newcommand{\diag}{\operatorname{diag}}
\newcommand{\codim}{\operatorname{codim}}
\newcommand{\avg}{\operatorname{avg}}
\newcommand{\im}{\operatorname{im}}
\newcommand{\ceil}[1]{\left\lceil #1\right\rceil}
\newcommand{\floor}[1]{\left\lfloor #1\right\rfloor}
\newcommand{\abs}[1]{\left| #1\right|}
\newcommand{\epsstar}{\varepsilon^*}
\newcommand{\one}{\mathbf 1}
\newcommand{\E}{\mathbb E}
\newcommand{\wt}{\operatorname{wt}}
\newcommand{\agr}{\operatorname{agr}}
\newcommand{\Rone}{\mathrm{R1}}
\newcommand{\Cen}{\operatorname{Cen}}

\crefname{conjecture}{Conjecture}{Conjectures}
\Crefname{conjecture}{Conjecture}{Conjectures}
\crefname{problem}{Problem}{Problems}
\Crefname{problem}{Problem}{Problems}
\crefname{residualinput}{Residual Input}{Residual Inputs}
\Crefname{residualinput}{Residual Input}{Residual Inputs}

\title{The Reed--Solomon MCA Frontier for the Proximity Prize\\[0.35em]
\textnormal{\normalsize A Self-Contained Conditional Closure Program, Exact Unsafe Certificates,\\
and the Finite Adjacent-Pair Conjecture}}
\author{Przemek Chojecki\\\small ulam.ai}
\date{Companion condensed manuscript matched to the main raw-working source, July 2026}

\begin{document}
\maketitle

\begin{abstract}
We isolate the Reed--Solomon mutual correlated agreement (MCA) problem that controls the Proximity Prize and write the current frontier as a single mathematical paper rather than as a research notebook.  This is a focused companion to the main raw-working manuscript \emph{Universal Field-Size Caps and a Two-Sided Sandwich for Mutual Correlated Agreement on Smooth Reed--Solomon Domains}.  The paper has two theorem-level sides and one deliberately exposed conditional side, following the main raw-working convention that the final frontier is a Conjecture, the three surviving near-capacity proof obligations are Residual Inputs, and the remaining auxiliary tasks are Questions.

The unsafe side is unconditional.  For a Reed--Solomon row with evaluation domain contained in a base field \(\BB\), split locators with a prescribed leading prefix give large lists at one received word.  At identity scale this produces the exact floor
\[
        |\List(\RS[\F,D,K],1-m/n,U)|\ge \frac{\binom nm}{|\BB|^{m-K}},
\]
up to the necessary integer ceiling.  A simple-pole conversion turns such a list for \(K=k+1\) into many bad MCA slopes for \(K=k\), and the flexible-budget version of the conversion pays only the cryptographic slope budget rather than the older \(q/k\) threshold.  This proves the deployed unsafe agreements \(1116047\) for the KoalaBear sextic MCA row and \(1116023\) for the Mersenne-31 circle line-round MCA row; the corresponding list-row unsafe agreements are \(1116046\) and \(1116022\).  All four comparisons are exact integer certificates.

The unconditional safe side is also made self-contained.  We prove a deep-regime MCA theorem below one third of the minimum distance; prove that, below half the minimum distance, MCA is bounded by ordinary correlated agreement plus at most the tangent-slope term; and include an elementary half-Johnson correlated-agreement bound derived only from the interleaved Johnson list bound.  Importing the standard unique-decoding proximity-gap theorem of Ben-Sasson--Carmon--Ishai--Kopparty--Saraf moves the generic safe edge to half distance.

The remaining near-capacity safe side is reduced to three explicit inputs: \(Q\), prefix/quotient flatness; \(BC\), a base-field-normalized split-pencil census; and \(SP\), primitive shift-pair control.  The reductions are included: interpolation-function slope elimination, rank-one support census, two-generator lattice dichotomy, split-pencil determinantal form, moment-kernel and collision-gap formulations of \(Q\), and the finite numerator staircase.  Thus the exact Proximity Prize prediction becomes a falsifiable adjacent-pair theorem: prove the finite versions of \(Q,BC,SP\) with constants fitting the printed margins, and the first safe agreements are exactly one more than the exact unsafe agreements above.  Prove the asymptotic forms with polynomial or \(e^{o(n)}\) losses, and the Reed--Solomon MCA threshold is
\[
        \delta_C^*(\varepsilon^*)=1-\rho-g^*(\rho,\log_2|\BB|)+o(1),
        \qquad
        H_2(\rho+g^*)=(\log_2|\BB|)g^* .
\]
The manuscript is therefore a settlement program in the strict sense: every unconditional ingredient is proved, every imported statement is isolated, and every remaining proof obligation is named as a concrete finite or asymptotic Residual Input.
\end{abstract}

\medskip
\noindent\textbf{Keywords:} Reed--Solomon codes, mutual correlated agreement, proximity gaps, Proximity Prize, prefix fibers, list decoding, FRI, STARKs.

{\small\setcounter{tocdepth}{2}\tableofcontents}


\section*{Reader's guide: theorem map and status}
\addcontentsline{toc}{section}{Reader's guide: theorem map and status}

The purpose of this version is not to hide the conditional part behind notation.  The paper is organized so that a reader can check exactly which statements are proved, which are imported, and which are the remaining settlement inputs.  The following table is the dependency map used throughout.

\begin{center}\small
\begin{tabular}{@{}p{0.22\linewidth}p{0.30\linewidth}p{0.36\linewidth}@{}}
\toprule
Component & Status & Role in the frontier \\
\midrule
Simple-pole conversion, prefix floors, exact deployed unsafe certificates & theorem-level, elementary, exact integer arithmetic & proves the last unsafe agreement on the four deployed MCA/list rows. \\
Deep-regime MCA and MCA-from-CA pincer & theorem-level, elementary, all linear codes & supplies unconditional safe anchors and removes the support-mismatch layer up to half distance. \\
Unique-decoding CA proximity gap & imported from \cite{BCIKS20} & moves the safe anchor for MCA to half the distance; the paper records exactly where the import is used. \\
Input Q & residual finite/asymptotic counting theorem & says prefix and quotient fibers are flat up to the allowed loss once quotient-periodic structure is paid. \\
Input BC & residual finite/asymptotic counting theorem & bounds the primitive split-pencil census left after slope elimination and lattice reduction. \\
Input SP & residual finite/asymptotic counting theorem & controls primitive shift-pair strata in the second-moment and collision ledger. \\
Active closure interface & main raw-working status convention & says broad-band, rank-one, balanced-core, and split-pencil statements are reduction stages; the final closure consumes only Q, BC, and SP. \\
Integer staircase & theorem-level bookkeeping & converts one unsafe integer and one safe integer into the exact closed-ball threshold. \\
\bottomrule
\end{tabular}
\end{center}

The phrase ``settles RS MCA'' is therefore used in a precise conditional sense.  The manuscript proves the negative side and the generic safe anchors; it reduces the remaining positive side to Q, BC, and SP.  A proof of their asymptotic forms gives the advertised asymptotic RS MCA frontier.  A proof of their finite forms with the constants demanded in \cref{tab:cert-margins} gives the exact Proximity Prize adjacent pairs in \cref{conj:finite-frontier}.  Conversely, any failure of the adjacent-pair prediction must appear as a concrete counterexample to one of Q, BC, or SP, not as an unspecified gap in the argument.

\paragraph{Relationship with the main raw-working manuscript.}
This paper is the compact companion to the main raw-working file
\texttt{cap25\_cap\_v13\_raw.tex}.  It is not a competing statement of the
theorem.  The raw-working file is the extended master source: it
keeps the universal field-size cap, quotient ledgers, map-smooth and
Chebyshev/circle variants, genus-one first-step material, threshold compilers,
L1/M1 residue-line tools, calibration data, and the final active-input cleanup.
Relative to the previous compact source, the master raw-working file adds several
useful clarifications: the grand-challenge cap is explicitly framed as a coarse
first-tier negative band rather than the final frontier; the finite certificate
part begins with a status table separating proved certificates, imported safe
anchors, and residual inputs; the final safe package is said to consist of
exactly Q, BC, and SP; the quotient, rank-one, balanced-core, tangent, and
prefix-collision sections are described as reductions toward these inputs, not
extra parallel assumptions; and the finite moment-order calibration is made
row-dependent.  The present paper is the companion core: it keeps only the
material needed to state the RS MCA Proximity Prize frontier, prove the unsafe
side, record the unconditional safe anchors, and expose the same safe-side
inputs.

The dictionary between the main raw-working file and this companion is as follows.
\begin{center}\scriptsize
\begin{tabular}{@{}>{\raggedright\arraybackslash}p{0.30\linewidth}>{\raggedright\arraybackslash}p{0.24\linewidth}>{\raggedright\arraybackslash}p{0.40\linewidth}@{}}
\toprule
Main raw-working manuscript & Matched companion manuscript & Meaning \\
\midrule
Conjecture \texttt{prob:capff1-frontier} & \cref{conj:asymptotic-frontier,conj:finite-frontier} and \cref{tab:cert-margins} & The same asymptotic entropy frontier and the same four deployed adjacent predictions. \\
Residual Input \texttt{prob:capg-active-Q} & Q, \cref{prob:Q} & Quotient and prefix flatness, with separate asymptotic and finite grades. \\
Residual Input \texttt{prob:capg-active-BC} & BC, \cref{prob:BC} & Base-field-normalized split-pencil census for interior balanced profiles. \\
Residual Input \texttt{prob:capg-active-shiftpairs} & SP, \cref{prob:SP} & Primitive shift-pair control after quotient-pullback pairs are deleted. \\
Active interface\newline\texttt{subsec:capg-active-}\newline\texttt{interface};\newline\texttt{prop:capg-active-}\newline\texttt{interface} & Reader's guide and \cref{thm:frontier-synthesis,thm:conditional-frontier} & Intermediate aperiodic, quotient, rank-one, and split-pencil statements are proof waypoints; only Q, BC, and SP are consumed by the final closure. \\
Proposition \texttt{prop:capg-final-active-}\newline\texttt{package} & \cref{thm:conditional-frontier} & The same conditional closure: polynomial losses give the asymptotic frontier with reserve; exact constants are needed for finite adjacent pairs. \\
Corollary \texttt{cor:capg-adjacent-pairs} & \cref{conj:finite-frontier} and \cref{tab:cert-margins} & The same last-unsafe/first-safe integer rows and margins. \\
Split-pencil active readings\newline\texttt{prob:capfr1-}\newline\texttt{normalized-band};\newline\texttt{rem:capfr1-split-}\newline\texttt{pencil-active-reading};\newline\texttt{rem:capfp-balanced-}\newline\texttt{active-reading} & \cref{sec:splitpencil,sec:Q} & The raw-working file marks the balanced-core and split-pencil forms as readings of the proof route, with base-field normalization delegated to BC. \\
Risk register and audit\newline\texttt{subsec:capfr1-risk};\newline\texttt{prob:capfr1-rung-audit} & \cref{sec:finite-obligations,app:cert} & Finite constants must be audited before claiming the adjacent-pair theorem, especially for the tight Mersenne-31 rows. \\
Questions \texttt{prob:explicit},\newline\texttt{prob:errorone}, \texttt{prob:isogeny} & Appendix discussion and boundary notes & Auxiliary directions only; they are not hidden assumptions in the frontier theorem. \\
Coarse-cap placement\newline\texttt{rem:twotier} & Reader's guide and introduction & The grand-challenge cap is a first-tier negative band; the active frontier is the later identity-scale/flexible-budget package. \\
Finite moment-order calibration\newline\texttt{rem:capfp-finite-orders} & \cref{rem:finite-moment-orders} & The required moment depth for a direct finite proof is row-dependent; the tight Mersenne-31 rows demand a substantially deeper hierarchy than a generic order-\(10^5\) slogan. \\
Conversion-threshold question\newline\texttt{ques:conversion-threshold} & Appendix boundary notes & Improving the simple-pole/flexible-budget threshold is a further direction, not an assumption in the current frontier package. \\
\bottomrule
\end{tabular}
\end{center}

Thus the two files are consistent precisely when the short paper is read as a
condensed companion to the main raw-working source, not as a replacement for every
auxiliary result in that source.  Any result present only in the main raw-working
file can be cited from there; any claim about the final Proximity Prize frontier
should use the status convention of this paper: unsafe side proved, generic
safe anchors proved/imported, final near-capacity safe side reduced to Q, BC,
and SP.  This is the companion counterpart of the main raw-working source's active interface: \texttt{subsec:capg-active-interface} and \texttt{prop:capg-active-interface}.

\begin{remark}[alignment with the main raw-working deltas]
The present companion incorporates the substantive differences between the
previous compact source and the current master source.  The source-level delta
has no change to the four printed frontier integers.  It changes the exposition
around \texttt{rem:twotier}, inserts the finite-certificate status table, states
that there are exactly three final residual inputs, replaces the fixed
``order \(10^5\)'' finite-moment slogan by the row-dependent calibration of
\cref{rem:finite-moment-orders}, and labels the conversion-threshold question.
\end{remark}

\begin{theorem}[frontier closure, informal synthesis]\label[theorem]{thm:frontier-synthesis}
For smooth multiplicative Reed--Solomon rows and the circle line-round rows considered here, the following statements hold.
\begin{enumerate}[label=\textup{(\alph*)}]
\item The unsafe side at the deployed rows is proved exactly by prefix floors and the flexible simple-pole conversion.
\item The safe side below one third of the distance is proved for MCA without imports; below half distance MCA is reduced to CA plus a tangent term.
\item Assuming Q, BC, and SP in the asymptotic polynomial-loss form, the RS MCA threshold satisfies
\[
        \delta_C^*(\varepsilon^*)=1-\rho-g^*(\rho,\log_2|\BB|)+o(1)
\]
with logarithmic agreement reserve.
\item Assuming Q, BC, and SP in the finite exact-constant form at the four deployed adjacent rows, \cref{conj:finite-frontier} holds.
\end{enumerate}
\end{theorem}

\begin{proof}
This is a guide to the later formal statements.  Item \textup{(a)} is \cref{thm:exact-unsafe}.  Item \textup{(b)} is \cref{thm:deep-mca,thm:mca-from-ca,thm:elementary-ca,cor:half-import}.  Items \textup{(c)} and \textup{(d)} are the two clauses of \cref{thm:conditional-frontier}, after the integer localization theorem \cref{thm:corridor}.  The reductions showing that Q, BC, and SP are the only remaining cells are given in \cref{sec:splitpencil,sec:Q,sec:finite-obligations}.
\end{proof}


\section{The problem and the claimed frontier}\label{sec:intro}

Let \(C=\RS[\F,D,k]\subseteq\F^n\) be the Reed--Solomon code obtained by evaluating polynomials of degree \(<k\) on a set \(D\) of size \(n\), and write \(\rho=k/n\).  In proximity testing one samples a random linear combination \(f_1+\gamma f_2\).  The mutual correlated agreement question asks whether, whenever this combination is close to \(C\), the two words \(f_1,f_2\) are already explained by codewords on the same large agreement set.  The Proximity Prize asks for sharp thresholds for this phenomenon in the Reed--Solomon rows used by modern proof systems.

There are two sharply different parts of the story.

First, the unsafe side is now elementary and essentially optimal for the present method.  Choose many subsets \(M\subseteq D\) of size \(m>k\).  The split locator
\[
        \Lambda_M(X)=\prod_{x\in M}(X-x)
\]
starts with \(w=m-k\) free leading coefficients beyond the code degree.  Pigeonholing these prefixes over the base field \(\BB\) gives a received word with at least \(\binom nm/|\BB|^w\) nearby codewords.  A simple-pole conversion turns this list into bad correlated-agreement slopes.  Optimizing the inequality
\[
        \binom nm\gtrsim |\BB|^w\cdot \epsstar |
        \F|
\]
leads to the entropy equation \(\Hb(\rho+g)=\beta g\), where \(m=(\rho+g)n\) and \(\beta=\log_2|\BB|\).

Second, the safe side is the hard part.  The safe-side theorem must say that every other source of bad slopes is smaller than the prize budget beyond the same entropy frontier.  The present paper reduces that theorem to three named inputs: prefix/quotient flatness, a split-pencil census, and primitive shift-pair control.  These inputs are not hidden inside the proof; they are stated explicitly in \cref{sec:conditional}.  The conditional theorem then follows by a finite integer staircase argument.

\subsection{The finite deployed prediction}

The deployed rows relevant here have
\[
        n=2^{21},\qquad k=2^{20},\qquad \rho=\frac12 .
\]
For the KoalaBear row, \(p=2^{31}-2^{24}+1\), \(\BB=\F_p\), \(\F=\F_{p^6}\), and the target is \(2^{-128}\).  For the Mersenne-31 circle line-round row, \(p'=2^{31}-1\), \(\BB=\F_{p'}\), \(\F=\F_{p'^4}\), and the target is \(2^{-100}\).  The circle row is Reed--Solomon after the standard twin-coset/torus change of variables; \cref{sec:circle} recalls the needed form.

The active finite prediction is the following adjacent-pair statement.

\begin{conjecture}[finite Proximity Prize frontier]\label[conjecture]{conj:finite-frontier}
At the deployed rows above, the first safe agreement is one more than the last exact unsafe agreement in the table:
\begin{center}
\scriptsize
\begin{tabular}{@{}lccc@{}}
\toprule
row & proved unsafe \(a_0\) & first safe \(a_0+1\) & fail margin \\
\midrule
KoalaBear MCA, target \(2^{-128}\) & \(1116047\) & \(1116048\) & \(22.2\) bits \\
KoalaBear list, target \(2^{-128}\) & \(1116046\) & \(1116047\) & \(22.0\) bits \\
Mersenne-31 MCA, target \(2^{-100}\) & \(1116023\) & \(1116024\) & \(3.3\) bits \\
Mersenne-31 list, target \(2^{-100}\) & \(1116022\) & \(1116023\) & \(3.1\) bits \\
\bottomrule
\end{tabular}
\end{center}
For an MCA row this is an adjacent closed-ball statement.  Unsafety is proved
at agreement threshold \(a_0\), i.e. throughout the printed near-capacity interval
starting at radius \((n-a_0)/n\).  Safety, if the finite Q/BC/SP certificates are
proved, begins at the next agreement \(a_0+1\), i.e. through radius
\((n-a_0-1)/n\).  Thus the displayed number \((n-a_0)/n\) is the first unsafe
radius, or equivalently the supremal boundary of the one-step safe interval;
the endpoint itself is not claimed safe under closed Hamming balls.
\end{conjecture}

Thus the two MCA endpoints predicted by the finite conjecture are
\[
        \delta^*_{\mathrm{KB}}(2^{-128})=\frac{981105}{2097152}=0.467827320\ldots,
        \qquad
        \delta^*_{\mathrm{M31}}(2^{-100})=\frac{981129}{2097152}=0.467838764\ldots .
\]
The unsafe inequalities at \(a_0\) are theorems.  The safe inequalities at \(a_0+1\) are conditional on the three finite safe-side inputs of \cref{sec:conditional}.  The small Mersenne margins explain why a polynomial-loss theorem is not enough for the exact finite prize: at \(n=2^{21}\), even a factor \(n\) costs \(21\) bits.

\subsection{The asymptotic frontier}

Let \(\Hb(x)=-x\log_2x-(1-x)\log_2(1-x)\).  For a base-field bit size \(\beta\), define
\begin{equation}\label{eq:gstar}
        g^*(\rho,\beta):=\sup\{g\in(0,1-\rho]:\Hb(\rho+g)\ge \beta g\}.
\end{equation}
This is the point where the identity-prefix lower floor has the same entropy as the target denominator supplied by the base field.

\begin{conjecture}[asymptotic Reed--Solomon MCA frontier]\label[conjecture]{conj:asymptotic-frontier}
Let \(D\subseteq\BB\) be a smooth multiplicative row or a circle line-round row, let \(C=\RS[\F,D,k]\) have rate \(\rho\), and let the target \(\epsstar\) be in the deployed cryptographic range, so that \(\log_2(\epsstar)^{-1}=O(1)\) compared with \(n\).  Then
\[
        \delta^*_C(\epsstar)=1-\rho-g^*(\rho,\log_2|\BB|)+o(1).
\]
More precisely, the equality holds with any agreement reserve \(R\gg\log n\) once the three asymptotic safe-side inputs in \cref{sec:conditional} are proved.
\end{conjecture}

For \(\rho=1/2\), \(\log_2(2^{31}-2^{24}+1)=30.988685\ldots\), so the predicted KoalaBear endpoint is \(0.4678266\ldots\); for \(p'=2^{31}-1\), it is \(0.4678383\ldots\).  The finite unsafe MCA endpoints are within about \(10^{-6}\) of these entropy-frontier values.

\section{Definitions}\label{sec:def}

Throughout \(\BB\subseteq\F\) are finite fields, \(q=|\F|\), and \(D\subseteq\BB\) has size \(n\).  The Reed--Solomon code is
\[
        \RS[\F,D,k]=\{(P(x))_{x\in D}:P\in\F[X],\ \deg P<k\}\subseteq\F^D.
\]
For words \(u,v\in\F^D\), \(\dist(u,v)\) denotes relative Hamming distance.  For \(S\subseteq D\), \(\distS\) denotes the same distance restricted to \(S\).  For a code \(C\), set
\[
        \List(C,\delta,u)=\{c\in C:\dist(u,c)\le\delta\},
        \qquad
        \List(C,\delta)=\max_u|\List(C,\delta,u)|.
\]
For an interleaved pair, \(C^{\equiv2}=C\times C\) is equipped with column distance:
\[
        \dist_2\bigl((f_1,f_2),(c_1,c_2)\bigr)
        =\frac{1}{n}\abs{\{x:f_1(x)\ne c_1(x)\text{ or }f_2(x)\ne c_2(x)\}}.
\]

\begin{definition}[correlated agreement]\label[definition]{def:ca}
For \(0\le\delta_{\rm fld}\le\delta_{\rm int}<1\), define
\[
\eca(C,\delta_{\rm fld},\delta_{\rm int})
:=\max_{f_1,f_2\in\F^D}\Pr_{\gamma\leftarrow\F}
\left[
\begin{array}{l}
        \dist(f_1+\gamma f_2,C)\le\delta_{\rm fld},\\
        \dist_2((f_1,f_2),C^{\equiv2})>\delta_{\rm int}
\end{array}\right].
\]
We write \(\eca(C,\delta)=\eca(C,\delta,\delta)\).
\end{definition}

\begin{definition}[mutual correlated agreement]\label[definition]{def:mca}
For \(0<\delta<1\), define
\[
\emca(C,\delta):=\max_{f_1,f_2\in\F^D}\Pr_{\gamma\leftarrow\F}
\left[
\begin{array}{l}
\exists S\subseteq D,\\ |S|\ge(1-\delta)n,\quad
\distS(f_1+\gamma f_2,C)=0,\\
\distS((f_1,f_2),C^{\equiv2})>0
\end{array}\right].
\]
A slope \(\gamma\) realizing the event is called MCA-bad for the pair.
\end{definition}

\begin{definition}[threshold]\label[definition]{def:threshold}
For a target \(\epsstar\in(0,1)\), set
\[
        \delta^*_C(\epsstar)=\sup\{\delta\in(0,1-\rho):\emca(C,\delta)\le\epsstar\},
        \qquad \rho=k/n.
\]
The convention is closed Hamming balls.  If agreement \(a\) means at least \(a\) matching positions, then a real radius \(\delta\) requires \(a(\delta)=\ceil{(1-\delta)n}\).
\end{definition}

\begin{fact}\label[fact]{fact:ca-to-mca}
For every linear code \(C\) and every \(\delta\),
\[
        \eca(C,\delta)\le\emca(C,\delta).
\]
Moreover \(\delta\mapsto\emca(C,\delta)\) is nondecreasing.
\end{fact}

\begin{proof}
If \(\gamma\) is CA-bad, choose a set \(S\) of size at least \((1-\delta)n\) on which \(f_1+\gamma f_2\) agrees with a codeword.  Since the pair \((f_1,f_2)\) is not \(\delta\)-close to \(C^{\equiv2}\), it cannot be simultaneously explained on this \(S\).  Thus the same \(\gamma\) is MCA-bad.  Monotonicity follows because the same witness set remains valid when the radius is increased.
\end{proof}

\section{The simple-pole conversion}\label{sec:conversion}

The core conversion takes a large list in the one-dimensional larger Reed--Solomon code \(C^+=\RS[\F,D,k+1]\) and turns it into many bad slopes for \(C=\RS[\F,D,k]\).  The idea is to divide by a simple pole \(x-\alpha\) with \(\alpha\notin D\).  A listed polynomial \(P\) then gives the slope \(P(\alpha)\):
\[
        \frac{U(x)}{x-\alpha}+P(\alpha)\left(-\frac{1}{x-\alpha}\right)
        =\frac{P(x)-P(\alpha)}{x-\alpha},
\]
which has degree \(<k\) on the agreement set.

\begin{theorem}[deep-point list-to-CA conversion]\label[theorem]{thm:deep-point}
Let \(C=\RS[\F,D,k]\), \(C^+=\RS[\F,D,k+1]\), \(q=|\F|>n=|D|\).  Let \(\delta\in(0,1)\), put \(f=\floor{\delta n}\), and assume
\[
        f\le n-k-1.
\]
For \(\eta\in[0,1)\), if
\[
        \eca(C,\delta)\le \eta\left(\frac{1}{k}-\frac{n}{kq}\right),
\]
then
\[
        \List(C^+,\delta)
        \le
        \ceil{ \frac{q\,\eca(C,\delta)}{1-
        \eta} } .
\]
Equivalently, if \(C^+\) has a ball of radius \(\delta\) with a list larger than the right-hand side, then \(\eca(C,\delta)\) exceeds the displayed threshold.
\end{theorem}

\begin{proof}
Let \(U:D\to\F\) be a received word and let \(P_1,\ldots,P_L\in\F[X]_{<k+1}\) be distinct listed polynomials.  Each agrees with \(U\) on at least
\[
        a:=n-f\ge k+1
\]
points of \(D\).  Put \(\Omega=\F\setminus D\).  For \(\alpha\in\Omega\), define
\[
        f_\alpha(x)=\frac{U(x)}{x-\alpha},
        \qquad
        g_\alpha(x)=-\frac{1}{x-\alpha}.
\]
If \(P_i\) agrees with \(U\) on \(S_i\), then with \(z_i(\alpha)=P_i(\alpha)\),
\[
        f_\alpha(x)+z_i(\alpha)g_\alpha(x)
        =\frac{P_i(x)-P_i(\alpha)}{x-\alpha}
\]
on \(S_i\), and the right side is a polynomial of degree \(<k\).  Thus every distinct value among \(P_i(\alpha)\) gives a close slope.

The pair \((f_\alpha,g_\alpha)\) is column-far: no degree \(<k\) polynomial can agree with \(g_\alpha\) on more than \(k\) points, because otherwise \((X-\alpha)G(X)+1\) would have more than \(k\) roots while having degree at most \(k\) and value \(1\) at \(X=\alpha\).  Since \(a>k\), the pair is not \(\delta\)-close to \(C^{\equiv2}\).  Hence distinct values \(P_i(\alpha)\) give CA-bad slopes.

It remains to choose \(\alpha\) with many distinct values.  For \(i\ne j\), \(P_i-P_j\) has degree at most \(k\), hence has at most \(k\) roots in \(\Omega\).  If \(M(\alpha)\) is the number of distinct values among \(P_1(\alpha),\ldots,P_L(\alpha)\), then averaging over \(\Omega\) and applying Cauchy--Schwarz gives some \(\alpha\) with
\[
        M(\alpha)
        \ge
        \frac{L(q-n)}{q-n+kL}.
\]
Therefore
\[
        \eca(C,\delta)
        \ge \frac{1}{q}\frac{L(q-n)}{q-n+kL}.
\]
Solving for \(L\), using the hypothesis
\(\eca(C,\delta)\le \eta(q-n)/(kq)\), gives
\[
        L\le \frac{q\,\eca(C,\delta)}{1-
        \eta}.
\]
Taking the maximum over \(U\) and rounding proves the theorem.
\end{proof}

For the finite prize we need a sharper operating point: if the target budget is much smaller than \(q/k\), the proof above can be stopped before the Cauchy--Schwarz threshold specialization.

\begin{corollary}[flexible-budget conversion]\label[corollary]{cor:flexible}
Let \(C=\RS[\F,D,k]\), \(C^+=\RS[\F,D,k+1]\), and \(q=|\F|>n\).  Let \(m\) be an integer with \(k+1\le m\le n\).  Suppose some received word \(U\) has
\[
        |\List(C^+,1-m/n,U)|\ge L_0
        \qquad\text{and}\qquad
        \binom{L_0}{2}k<q-n.
\]
Then some simple-pole line has at least \(L_0\) distinct CA-bad slopes at agreement threshold \(m\).  Consequently
\[
        \eca(C,1-m/n)\ge \frac{L_0}{q},
        \qquad
        \emca(C,\delta)\ge \frac{L_0}{q}\quad\text{for every }\delta\ge1-m/n.
\]
If the pole is restricted to \(\F\setminus(D\cup\BB)\), the same conclusion holds with \(q-n\) replaced by \(q-|\BB|\), and the certifying line is not \(\BB\)-rational.
\end{corollary}

\begin{proof}
Take \(L_0\) listed polynomials.  For \(i\ne j\), \(P_i-P_j\) has at most \(k\) roots on \(\F\setminus D\).  Hence the total number of pair collisions over \(\Omega=\F\setminus D\) is at most \(k\binom{L_0}{2}<|\Omega|\), so some pole \(\alpha\in\Omega\) has no collisions at all.  At this pole the \(L_0\) values \(P_i(\alpha)\) are distinct, and the proof of \cref{thm:deep-point} shows that they are CA-bad slopes.  The MCA statement follows from \cref{fact:ca-to-mca}.  The restricted-pole variant averages over \(\F\setminus(D\cup\BB)\).
\end{proof}

\section{Prefix floors}\label{sec:floors}

\subsection{Identity scale}

The identity scale is the terminal scale of all quotient constructions.  It is also the scale that reaches the deployed entropy frontier.

\begin{lemma}[identity-prefix floor]\label[lemma]{lem:identity-floor}
Let \(\BB\subseteq\F\), let \(D\subseteq\BB\) have \(|D|=n\), and let \(1\le K\le n\).  For every integer \(m\) with \(K\le m\le n\), put \(w=m-K\).  Then there is a \(\BB\)-valued received word \(U:D\to\BB\) such that
\[
        |\List(\RS[\F,D,K],1-m/n,U)|
        \ge
        \ceil{ \frac{\binom nm}{|\BB|^w} }.
\]
\end{lemma}

\begin{proof}
For \(M\subseteq D\), \(|M|=m\), write
\[
        \Lambda_M(X)=\prod_{x\in M}(X-x)
        =X^m+\sum_{j=1}^{m}(-1)^j e_j(M)X^{m-j}.
\]
For \(z=(z_1,\ldots,z_w)\in\BB^w\), define
\[
        U_z=\bigl(X^m+\sum_{j=1}^w z_jX^{m-j}\bigr)\bigm|_D.
\]
The map \(M\mapsto((-1)^j e_j(M))_{1\le j\le w}\) has codomain \(\BB^w\), so one fiber has at least \(\binom nm/|\BB|^w\) members.  Fix a heaviest prefix \(z^*\).  For \(M\) in this fiber put
\[
        c_M=U_{z^*}-\Lambda_M|_D
        =-\sum_{j=w+1}^{m}(-1)^j e_j(M)X^{m-j}\bigm|_D.
\]
Since \(m-w-1=K-1\), \(c_M\in\RS[\F,D,K]\).  It agrees with \(U_{z^*}\) on the \(m\) roots of \(\Lambda_M\) in \(D\).  If two members \(M,M'\) give the same codeword, then the difference of the two displayed polynomials has degree at most \(K-1<n\) and vanishes on all of \(D\), hence is zero.  The tail coefficients and the fixed prefix then determine the full monic locator, so \(M=M'\).  Thus the listed codewords are distinct.
\end{proof}

The prefix construction is not merely a lower bound for a chosen heaviest word; every word with prescribed prefix has exactly the corresponding fiber as its list.

\begin{proposition}[exact witness lists]\label[proposition]{prop:exact-witness-list}
In the setting of \cref{lem:identity-floor}, for every \(z\in\BB^w\),
\[
        \List(\RS[\F,D,K],1-m/n,U_z)
        =\{c_M:M\subseteq D,\ |M|=m,\ ((-1)^j e_j(M))_{j\le w}=z\}.
\]
Moreover no codeword agrees with \(U_z\) on more than \(m\) positions.
\end{proposition}

\begin{proof}
The inclusion \(\supseteq\) is the construction above.  Conversely, suppose \(c\in\RS[\F,D,K]\) agrees with \(U_z\) on at least \(m\) points, and let \(\hat c\in\F[X]\) be its polynomial of degree \(<K\).  Then
\[
        P(X)=X^m+\sum_{j=1}^wz_jX^{m-j}-\hat c(X)
\]
is monic of degree \(m\), because \(\deg\hat c<K=m-w\).  It has at least \(m\) roots in \(D\), so it is the locator of its root set \(M\subseteq D\), \(|M|=m\).  Comparing the leading \(w\) coefficients puts \(M\) in the prefix fiber of \(z\), and comparing the remaining coefficients gives \(c=c_M\).  A nonzero polynomial of degree \(m\) cannot have more than \(m\) roots.
\end{proof}

\subsection{Polynomial map scales}

The same construction works at every scale where \(D\) is partitioned into complete fibers of a polynomial map.  This is the right form for multiplicative quotients \(X^c\) and for Chebyshev folds on circle domains.

\begin{definition}[map-smooth scale]\label[definition]{def:mapsmooth}
A polynomial \(\varphi\in\BB[X]\) of degree \(c\ge2\) is a complete uniform scale for \(D\subseteq\BB\) if every fiber of \(\varphi:D\to\varphi(D)\) has size exactly \(c\).  Put \(Q=\varphi(D)\), so \(|Q|=n/c\).
\end{definition}

\begin{proposition}[graded prefix floor at a map-smooth scale]\label[proposition]{prop:map-floor}
Assume \(\varphi\) is a complete uniform scale of degree \(c\) for \(D\).  Let \(1\le K\le n\), let \(m\) satisfy \(\ceil{K/c}\le m\le |Q|\), and put
\[
        w=m-\ceil{K/c}\ge0.
\]
Then there is a \(\BB\)-valued word \(U:D\to\BB\) such that
\[
        |\List(\RS[\F,D,K],1-mc/n,U)|
        \ge
        \ceil{ \frac{\binom{|Q|}{m}}{|\BB|^w} }.
\]
\end{proposition}

\begin{proof}
For \(M\subseteq Q\), \(|M|=m\), use the locator
\[
        \Lambda_M(X)=\prod_{b\in M}(\varphi(X)-b)
        =\sum_{j=0}^m(-1)^j e_j(M)\varphi(X)^{m-j}.
\]
It vanishes on exactly \(mc\) points of \(D\).  Pigeonhole the first \(w\) coefficients in \(\BB^w\), and subtract the locator from the corresponding prefix word.  The remaining degree is at most
\[
        c(m-w-1)=c(\ceil{K/c}-1)\le K-1.
\]
The same degree and coefficient comparison as in \cref{lem:identity-floor} proves injectivity.
\end{proof}

\subsection{A universal field-size cap}

The older universal cap is not the final deployed frontier, but it is important because it is completely field-size-uniform and elementary.

\begin{theorem}[universal unsafe cap]\label[theorem]{thm:universal-cap}
Let \(\BB\subseteq\F\), \(q=|\F|\), and let \(D\subseteq\BB^\times\) be a multiplicative coset of order \(n\).  Let \(k/n=\rho\), \(C=\RS[\F,D,k]\), and let \(N\mid n\) with \(a=n/N\mid k\) and \((1-\rho)N\ge3\).  If
\[
        \binom{N}{\rho N+2}\ge |\BB|\left(\frac{q}{k}+1\right),
\]
then, with \(\delta_N=1-\rho-2/N\),
\[
        \eca(C,\delta)>\frac{1}{2k}\left(1-\frac{n}{q}\right)
\]
for every \(\delta_N\le\delta\le1-\rho-1/n\), and
\[
        \emca(C,\delta)>\frac{1}{2k}\left(1-\frac{n}{q}\right)
\]
for every \(\delta_N\le\delta<1-\rho\).
\end{theorem}

\begin{proof}
Apply \cref{prop:map-floor} to \(\varphi(X)=X^a\), \(K=k+1\), and \(m=\rho N+2\).  It gives a \(C^+=\RS[\F,D,k+1]\) list of size at least \(\binom N{\rho N+2}/|\BB|\) at radius \(\delta_N\).  If \(\eca(C,\delta)\le(1/2k)(1-n/q)\) for some integer-admissible \(\delta\ge\delta_N\), \cref{thm:deep-point} with \(\eta=1/2\) gives
\[
        \List(C^+,\delta)<\frac{q}{k}+1,
\]
contradicting the assumed floor and list monotonicity.  The MCA statement follows from \cref{fact:ca-to-mca} at \(\delta_N\) and monotonicity.
\end{proof}

\begin{corollary}[grand-envelope cap]\label[corollary]{cor:grand-cap}
Let \(\rho\in\{1/2,1/4,1/8,1/16\}\), \(q<2^{256}\), and \(k=\rho n\le2^{40}\).  For smooth multiplicative rows with \(2^{10}\mid n\) at rates \(1/2,1/4,1/8\), and \(2^{11}\mid n\) at rate \(1/16\),
\[
        \emca(C,\delta)>2^{-86}>2^{-128}
\]
throughout
\[
\delta\in[1-\rho-2^{-9},1-\rho)\quad(\rho=1/2,1/4,1/8),
\]
respectively
\[
\delta\in[1-\rho-2^{-10},1-\rho)\quad(\rho=1/16).
\]
If \(q\ge2n\), the lower bound improves to \(2^{-42}\).
\end{corollary}

\begin{proof}
Use \(N=1024\) for \(\rho=1/2,1/4,1/8\) and \(N=2048\) for \(\rho=1/16\).  The binomial floor beats the worst-case denominator \(|\BB|(q/k+1)\) in the challenge envelope.  Since \(k\le2^{40}\) and \(q>n\),
\[
        \frac{1}{2k}\left(1-\frac{n}{q}\right)>2^{-86}
\]
under the stated rows, and if \(q\ge2n\) it is at least \(1/(4k)\ge2^{-42}\).
\end{proof}

\section{Exact deployed unsafe certificates}\label{sec:deployed-unsafe}

We now specialize the identity floor to the deployed rows.  The list rows use \(K=k\).  The MCA rows use \(K=k+1\) and the flexible-budget conversion.

\begin{theorem}[exact identity-scale unsafe certificates]\label[theorem]{thm:exact-unsafe}
Let \(n=2^{21}\), \(k=2^{20}\).
\begin{enumerate}[label=\textup{(\alph*)}]
\item \textup{KoalaBear MCA.}  Let \(p=2^{31}-2^{24}+1\), \(q=p^6\), and \(L_0=\floor{2^{-128}q}+1\).  With \(m=1116047\), \(K=k+1\), and \(w=m-K=67470\),
\[
        \binom n{1116047}>p^{67470}\floor{p^6/2^{128}},
        \qquad
        \binom{L_0}{2}k<q-n.
\]
Hence \(\emca(C,\delta)>2^{-128}\) for every
\[
        \delta\in\left[\frac{981105}{2097152},\frac{1}{2}\right).
\]
The same prefix comparison fails at \(m+1=1116048\) with \(w+1\).

\item \textup{KoalaBear list.}  With \(m=1116046\), \(K=k\), and \(w=67470\),
\[
        \binom n{1116046}\,2^{128}>p^{67470}p^6,
\]
and the same comparison fails at \(m+1=1116047\).  Thus the list-row unsafe edge is \(981106/2097152=490553/1048576\).

\item \textup{Mersenne-31 MCA.}  Let \(p'=2^{31}-1\), \(q=(p')^4\), and \(L_0=\floor{2^{-100}q}+1\).  With \(m=1116023\), \(K=k+1\), and \(w=67446\),
\[
        \binom n{1116023}>(p')^{67446}\floor{(p')^4/2^{100}},
        \qquad
        \binom{L_0}{2}k<q-n.
\]
Hence \(\emca(C,\delta)>2^{-100}\) for every
\[
        \delta\in\left[\frac{981129}{2097152},\frac{1}{2}\right).
\]
The comparison fails at \(m+1=1116024\) with \(w+1\).

\item \textup{Mersenne-31 list.}  With \(m=1116022\), \(K=k\), and \(w=67446\),
\[
        \binom n{1116022}\,2^{100}>(p')^{67446}(p')^4,
\]
and the same comparison fails at \(m+1=1116023\).  The list-row unsafe edge is \(981130/2097152=490565/1048576\).
\end{enumerate}
\end{theorem}

\begin{proof}
For list rows, \cref{lem:identity-floor} gives a list of size at least \(\ceil{\binom nm/p^w}\).  The displayed inequalities are exactly the statement that this list exceeds \(2^{-t}q\), with \(t=128\) or \(100\).  For MCA rows, \cref{lem:identity-floor} gives a \(C^+\)-list of size at least \(L_0\), and \cref{cor:flexible} converts it to \(L_0\) bad slopes.  Since \(L_0/q>2^{-t}\), the MCA target is exceeded.  All inequalities are exact integer comparisons; the approximate pass/fail margins are recorded in \cref{tab:cert-margins}.
\end{proof}

\begin{table}[t]
\centering
\begin{tabular}{@{}lrrrr@{}}
\toprule
row & \(m\) & \(w\) & pass margin & next-step margin \\
\midrule
KoalaBear MCA & 1116047 & 67470 & \(+8.98\) bits & \(-22.20\) bits \\
KoalaBear list & 1116046 & 67470 & \(+9.16\) bits & \(-22.01\) bits \\
Mersenne-31 MCA & 1116023 & 67446 & \(+27.93\) bits & \(-3.26\) bits \\
Mersenne-31 list & 1116022 & 67446 & \(+28.11\) bits & \(-3.07\) bits \\
\bottomrule
\end{tabular}
\caption{Logarithmic orientation margins for the exact inequalities of \cref{thm:exact-unsafe}.  Positive means the displayed unsafe comparison passes; negative means the same comparison at the next agreement fails.  The proof uses exact integer arithmetic, not these rounded margins.}
\label{tab:cert-margins}
\end{table}


\section{Field of definition and extension lines}\label{sec:extension}

The deployed KoalaBear row is an extension-field row: the domain lies in \(\BB=\F_p\), while slopes are sampled in \(\F=\F_{p^6}\).  It is therefore important that the accounting is not secretly one-dimensional over the wrong field.  The following restriction-of-scalars statement is the clean form.

Fix a \(\BB\)-basis \(1=\theta_0,\theta_1,\ldots,\theta_{e-1}\) of \(\F\), \(e=[\F:\BB]\).  Write any word \(w:D\to\F\) as \(w=\sum_i\theta_iw_i\), with \(w_i:D\to\BB\).

\begin{theorem}[restriction of scalars for Reed--Solomon rows]\label[theorem]{thm:restriction-scalars}
Let \(D\subseteq\BB\).  A word \(c:D\to\F\) lies in \(\RS[\F,D,k]\) if and only if each component \(c_i\) lies in \(\RS[\BB,D,k]\).  Consequently
\[
        \dist(w,\RS[\F,D,k])
        =\dist_e((w_0,\ldots,w_{e-1}),\RS[\BB,D,k]^{\equiv e}),
\]
where the right-hand side is column distance for the \(e\)-interleaved base-field code.  A support \(S\subseteq D\) explains \(w\) over \(\F\) if and only if it column-explains all components over \(\BB\).
\end{theorem}

\begin{proof}
If \(c\) is the evaluation of \(P\in\F[X]\), write \(P=\sum_i\theta_iP_i\) with \(P_i\in\BB[X]\); since \(D\subseteq\BB\), evaluation commutes with the decomposition and \(c_i=P_i|_D\).  The converse is the same argument in reverse.  Pointwise equality over \(\F\) is simultaneous equality of all \(\BB\)-components, which proves the distance and support statements.
\end{proof}

For a sampled slope \(\gamma\in\F\), multiplication by \(\gamma\) is a \(\BB\)-linear map on components.  Thus a received line \(f_1+\gamma f_2\) over \(\F\) is an \(e\)-interleaved \(\BB\)-linear sampler.  Every extension-valued bad slope is therefore visible in a base-field interleaved ledger.  Conversely, the simple-pole conversion of \cref{cor:flexible} can be run with poles \(\alpha\in\F\setminus\BB\), producing genuinely extension-valued certifying lines whenever the list is large enough.

\begin{corollary}[base-field confinement of base-valued lines]\label[corollary]{cor:confinement-core}
If \(f_1,f_2\) are \(\BB\)-valued words and \(\gamma\notin\BB\), then any support-wise explanation of \(f_1+\gamma f_2\) over \(\F\) splits into separate \(\BB\)-explanations of \(f_1\) and \(f_2\) on the same support.  Hence a \(\BB\)-valued received line has no genuinely extension-valued MCA-bad slopes.
\end{corollary}

\begin{proof}
Write \(\F=\BB\oplus V\) with \(1\) in the first component.  Since \(f_1,f_2\) are \(\BB\)-valued and \(\gamma\notin\BB\), the components of \(f_1+\gamma f_2\) along \(1\) and along the non-base direction generated by \(\gamma\) recover \(f_1\) and \(f_2\) by an invertible \(\BB\)-linear transformation.  By \cref{thm:restriction-scalars}, any codeword explanation over \(\F\) gives component codeword explanations over \(\BB\) on the same support.  Therefore the mutual condition cannot fail.
\end{proof}

This explains a useful distinction.  The deployed extension-field failures certified by the prefix floor are fiber-borne and genuinely extension-valued; they are obtained by choosing an extension pole, not by taking a base-valued adversarial line.


\section{Safe-side theorems that are already unconditional}\label{sec:safe}

The unconditional safe side has two elementary pieces.  First, below one third of the minimum distance, MCA has error at most \((r+1)/q\).  Second, below half the distance, MCA is controlled by ordinary correlated agreement up to an explicit tangent term.  The half-distance endpoint therefore follows from any CA theorem in the unique-decoding range; using the imported theorem of Ben-Sasson--Carmon--Ishai--Kopparty--Saraf gives the strongest current safe edge.

\begin{theorem}[deep-regime MCA]\label[theorem]{thm:deep-mca}
Let \(C\subseteq\F^n\) be a linear code of minimum Hamming weight \(w_{\min}\), let \(q=|\F|\), let \(\delta\in(0,1)\), and put \(r=\floor{\delta n}\).  If
\[
        3r\le w_{\min}-1,
\]
then
\[
        \eca(C,\delta)\le\frac{r+1}{q},
        \qquad
        \emca(C,\delta)\le\frac{r+1}{q}.
\]
For \(C=\RS[\F,D,k]\), the hypothesis is \(3r\le n-k\).
\end{theorem}

\begin{proof}
Fix \((f_1,f_2)\).  Let \(S_{\rm cl}\) be the set of slopes \(\gamma\) for which \(f_1+\gamma f_2\) is within \(r\) errors of \(C\).  If \(|S_{\rm cl}|\le1\), there is nothing to prove.  Otherwise choose two close slopes \(\gamma_1\ne\gamma_2\) with explaining codewords \(p_1,p_2\) and error sets \(E_1,E_2\), \(|E_i|\le r\).  Put
\[
        c_2=\frac{p_1-p_2}{\gamma_1-\gamma_2},
        \qquad
        c_1=p_1-\gamma_1c_2,
        \qquad
        e_i=f_i-c_i.
\]
Off \(T=E_1\cup E_2\), \(|T|\le2r\), the pair \((f_1,f_2)\) agrees with \((c_1,c_2)\).  For any other close slope \(\gamma\) with codeword \(p\) and error set \(E\), the codeword \(p-(c_1+
\gamma c_2)\) vanishes off \(T\cup E\), so has weight at most \(3r\le w_{\min}-1\).  It is zero.  Thus for every close slope, \(e_1+\gamma e_2\) has support size at most \(r\).

Let \(T'=\Supp(e_1,e_2)\subseteq T\).  If \(|T'|\ge r+1\), each close slope must make at least \(|T'|-r\) of the nonzero affine functions \(e_{1,j}+\gamma e_{2,j}\) vanish; each coordinate contributes at most one slope.  Double counting gives
\[
        |S_{\rm cl}|(|T'|-r)\le |T'|,
        \qquad\text{so}\qquad
        |S_{\rm cl}|\le r+1.
\]
If \(|T'|\le r\), the pair is tangent to \((c_1,c_2)\) outside at most \(r\) coordinates.  It is then not CA-bad at radius \(\delta\), and its MCA-bad slopes are among the at most \(|T'|\) tangent ratios \(-e_{1,j}/e_{2,j}\).  Hence every pair contributes at most \(r+1\) CA- or MCA-bad slopes.  Dividing by \(q\) gives the bounds.
\end{proof}

\begin{proposition}[exact tangent cell]\label[proposition]{prop:tangent-cell}
For \(C=\RS[\F,D,k]\), if \(A=n-r\) and \(3r\le n-k\), then the maximum number of finite MCA-bad slopes at agreement threshold \(A\) is exactly \(r+1\) for support-wise MCA/CA numerator accounting up to the harmless endpoint distinction; the upper bound is \cref{thm:deep-mca}, and matching examples are supported on \(r+1\) tangent coordinates.
\end{proposition}

\begin{proof}
The upper bound is \cref{thm:deep-mca}.  For the lower construction, choose \(T=\{t_0,\ldots,t_r\}\subseteq D\) and distinct slopes \(\gamma_0,
\ldots,\gamma_r\).  Let \(f_2(t_i)=1\), \(f_1(t_i)=-\gamma_i\), and let both words vanish off \(T\).  At slope \(\gamma_i\), the combination vanishes at \(t_i\) and off \(T\), hence agrees with the zero codeword on \(D\setminus(T\setminus\{t_i\})\), a set of size \(n-r=A\).  The inequalities implied by \(3r\le n-k\) ensure that no two degree \(<k\) explanations can mimic this on the witness support unless they are zero, so the pair is not simultaneously explained on that support.  Thus the \(r+1\) slopes are bad.
\end{proof}

\begin{theorem}[MCA from CA up to half distance]\label[theorem]{thm:mca-from-ca}
Let \(C\subseteq\F^n\) be a linear code of minimum Hamming weight \(w_{\min}\), let \(q=|\F|\), let \(\delta\in(0,1)\), and put \(r=\floor{\delta n}\).  If
\[
        2r\le w_{\min}-1,
\]
then
\[
        \emca(C,\delta)\le\max\left(\eca(C,\delta),\frac{r}{q}\right).
\]
For \(C=\RS[\F,D,k]\), this is the range \(2r\le n-k\).
\end{theorem}

\begin{proof}
Fix a pair \((f_1,f_2)\).  If it is farther than \(r\) columns from \(C^{\equiv2}\), then every MCA-bad slope has a witness set of size at least \(n-r\), hence the combination is \(r\)-close to \(C\); such a slope is CA-bad by definition.  If the pair is within \(r\) columns of \((p_1,p_2)\in C^{\equiv2}\), write \(e_i=f_i-p_i\), supported on a set \(E\) of size at most \(r\).  Let \(\gamma\) be MCA-bad, with witness set \(S\) and codeword \(c\).  On \(S\setminus E\), the codeword \(c-(p_1+\gamma p_2)\) vanishes; since \(|S\setminus E|\ge n-2r\), it has weight at most \(2r\le w_{\min}-1\), so it is zero.  Thus badness requires some \(j\in S\cap E\) with \(e_{1,j}+\gamma e_{2,j}=0\), and each \(j\in E\) gives at most one slope.  There are at most \(r\) such slopes.
\end{proof}

\begin{corollary}[conditional half-distance safe edge]\label[corollary]{cor:half-import}
Import the unique-decoding proximity-gap theorem of \cite{BCIKS20} in the form \(\eca(C,\delta)\le n/q\) for \(\delta\le(1-\rho)/2\).  Then for \(C=\RS[\F,D,k]\),
\[
        \emca(C,\delta)\le \frac{n}{q}
        \qquad\text{whenever}\qquad
        2\floor{\delta n}\le n-k.
\]
For the KoalaBear MCA row this gives safety at every \(\delta\le1/4\), since \(n/q<2^{-164}<2^{-128}\).  For the Mersenne-31 row it gives safety at every \(\delta\le1/4\), since \(n/q<2^{-102}<2^{-100}\).
\end{corollary}

\begin{proof}
Combine the import with \cref{thm:mca-from-ca}; \(r/q\le n/q\).  The numerical inequalities are immediate from the stated field sizes.
\end{proof}


\section{Self-contained list and CA inputs}\label{sec:johnson-ca}

The imported proximity-gap theorem is useful for the strongest safe anchor, but the paper also needs an import-free safe ledger.  This section records the elementary list-to-CA route.  It is not strong enough to close the near-capacity frontier, yet it is important because it verifies the proof grammar without hidden black boxes and it prices finite certificate cells that are outside the hard band.

\begin{theorem}[elementary Johnson list bound, uniform in interleaving arity]\label[theorem]{thm:johnson-list}
Let \(C=\RS[\F,D,k]\) on an evaluation set \(D\) of size \(n\).  Let \(s\ge1\), and equip \(C^{\equiv s}\) with column distance.  If \(a\) is an agreement threshold with
\[
        a^2>(k-1)n,
\]
then, for every received \(s\)-tuple \(U\),
\[
        \left|\List\left(C^{\equiv s},1-\frac{a}{n},U\right)\right|
        \le
        \frac{n(a-k+1)}{a^2-(k-1)n}.
\]
The same proof applies to any code family in which two distinct interleaved codewords column-agree on at most \(k-1\) coordinates.
\end{theorem}

\begin{proof}
Let \(\mathbf c^{(1)},\ldots,\mathbf c^{(L)}\) be distinct listed interleaved codewords, and let \(S_i\subseteq D\) be the column-agreement set of \(\mathbf c^{(i)}\) with \(U\).  Then \(|S_i|\ge a\).  If \(i\ne j\), the two interleaved codewords differ in at least one row, and two distinct degree \(<k\) polynomials agree on at most \(k-1\) points; hence
\[
        |S_i\cap S_j|\le k-1 .
\]
For \(x\in D\), let \(d(x)=|\{i:x\in S_i\}|\).  Then \(\sum_xd(x)\ge La\), and
\[
        \sum_x\binom{d(x)}2=\sum_{i<j}|S_i\cap S_j|
        \le \binom L2(k-1).
\]
By convexity of \(t(t-1)\),
\[
        \sum_x d(x)(d(x)-1)
        \ge n\left(\frac{La}{n}\right)\left(\frac{La}{n}-1\right),
\]
with the inequality vacuous if the right side is negative.  Combining the last two displays and cancelling one factor of \(L\) gives
\[
        a(La-n)\le (L-1)(k-1)n,
\]
which rearranges to
\[
        L\bigl(a^2-(k-1)n\bigr)
        \le n(a-k+1).
\]
The denominator is positive by hypothesis, proving the bound.
\end{proof}

\begin{theorem}[elementary correlated agreement at half the Johnson radius]\label[theorem]{thm:elementary-ca}
Let \(C=\RS[\F,D,k]\), \(|D|=n\), \(q=|\F|\).  Put \(r=\floor{\delta n}\).  Assume
\[
        n-2r>0,
        \qquad
        (n-2r)^2>(k-1)n.
\]
Define
\[
        L_2=\frac{n(n-2r-k+1)}{(n-2r)^2-(k-1)n}.
\]
Then
\[
        \eca(C,\delta)\le \frac{1+(r+1)L_2}{q}.
\]
More generally the same bound holds for \(\eca(C,\delta,\delta_{\rm int})\) for every \(\delta_{\rm int}\ge\delta\).
\end{theorem}

\begin{proof}
Fix a pair \((f_1,f_2)\) that is \(\delta_{\rm int}\)-far from \(C^{\equiv2}\), and let \(B\) be the set of CA-bad slopes.  For each \(\gamma\in B\), choose a codeword \(c_\gamma\in C\) and an agreement set \(S_\gamma\) with \(|S_\gamma|\ge n-r\) such that
\[
        f_1+\gamma f_2=c_\gamma\quad\text{on }S_\gamma.
\]
If \(|B|\le1\), there is nothing to prove.  Fix \(\gamma_0\in B\).  For any \(\gamma\in B\setminus\{\gamma_0\}\), set
\[
        P_2=\frac{c_{\gamma_0}-c_\gamma}{\gamma_0-\gamma},
        \qquad
        P_1=c_{\gamma_0}-\gamma_0P_2 .
\]
Then \((P_1,P_2)\in C^{\equiv2}\), and on \(S_{\gamma_0}\cap S_\gamma\), a set of size at least \(n-2r\), the two linear equations in \((f_1,f_2)\) force \((f_1,f_2)=(P_1,P_2)\).  Thus every \(\gamma\ne\gamma_0\) rides an interleaved codeword in
\[
        \mathcal L=
        \List\left(C^{\equiv2},1-\frac{n-2r}{n},(f_1,f_2)\right),
\]
whose size is at most \(L_2\) by \cref{thm:johnson-list}.

Now fix one rider pair \((P_1,P_2)\in\mathcal L\), and let \(E\) be the column error set on which \((f_1,f_2)\ne(P_1,P_2)\).  Since the original pair is \(\delta_{\rm int}\)-far and \(\delta_{\rm int}\ge\delta\), we have \(|E|\ge r+1\).  If a slope \(\gamma\) rides this pair, then
\[
        (f_1-P_1)+\gamma(f_2-P_2)
\]
vanishes on at least \(|E|-r\) points of \(E\), because \(f_1+\gamma f_2\) agrees with \(P_1+\gamma P_2\) on at least \(n-r\) coordinates.  Each coordinate of \(E\) contributes at most one slope to this vanishing condition.  Hence
\[
        \#\{\text{riders of }(P_1,P_2)\}\,(|E|-r)\le |E|,
\]
so the number of riders of this pair is at most \(|E|/(|E|-r)\le r+1\).  Summing over at most \(L_2\) rider pairs and adding back \(\gamma_0\) gives \(|B|\le1+(r+1)L_2\).  Divide by \(q\) and maximize over pairs.
\end{proof}

\begin{corollary}[import-free MCA safe cells]\label[corollary]{cor:self-contained-safe}
Under the hypotheses of \cref{thm:elementary-ca}, if in addition \(2r\le n-k\), then
\[
        \emca(C,\delta)\le \frac{1+(r+1)L_2}{q}.
\]
In the challenge rates this half-Johnson cell improves the elementary one-third-distance cell only below rate \(1/4\); at rates \(1/2\) and \(1/4\) the deep-regime theorem remains the stronger import-free MCA anchor.
\end{corollary}

\begin{proof}
Apply \cref{thm:mca-from-ca}.  Its tangent term is \(r/q\), while \(1+(r+1)L_2\ge r\) in the nontrivial cases, so the displayed bound dominates the maximum.
\end{proof}

\begin{remark}[why this section does not settle the frontier]\label[remark]{rem:johnson-limit}
The Johnson double count is a list theorem.  In the frontier window, the agreement is close to \(\rho n+g n\), much below the Johnson agreement \(\sqrt{(k-1)n}\) at rate \(1/2\).  Thus \cref{thm:elementary-ca} is a clean import-free proof cell, but not the near-capacity safe theorem.  The latter is exactly the Q/BC/SP package of \cref{sec:conditional,sec:splitpencil,sec:Q}.
\end{remark}

\section{Integer threshold compilers}\label{sec:staircase}

The finite prize is an integer problem.  The following bookkeeping is simple but prevents off-by-one mistakes.

\begin{definition}[integer numerator staircase]\label[definition]{def:staircase}
Let \(I=\{a_{\min},\ldots,a_{\max}\}\) be an agreement interval, and let \(N:I\to\mathbb Z_{\ge0}\) be nonincreasing.  For MCA, \(N(a)\) is the number of bad sampled slopes at agreement at least \(a\); for list rows, it is the list numerator.  Given a denominator \(Q\) and target \(\epsstar\), set
\[
        B_*:=\floor{\epsstar Q}.
\]
An agreement \(a\) is safe if \(N(a)\le B_*\) and unsafe if \(N(a)>B_*\).
\end{definition}

\begin{theorem}[staircase localization]\label[theorem]{thm:staircase}
For a nonincreasing numerator \(N\), exactly one of the following holds: all agreements are safe; all are unsafe; or there is a unique first safe agreement \(a_*\), with
\[
        N(a_*-1)>B_*\ge N(a_*).
\]
If this interior case holds, the real closed-ball safe supremum is
\[
        \frac{n-a_*+1}{n},
\]
and that endpoint is not attained.
\end{theorem}

\begin{proof}
The predicate \(N(a)\le B_*\) is monotone nondecreasing in \(a\).  This gives the unique first true point.  The real endpoint follows from \(a(\delta)=\ceil{(1-\delta)n}\): the condition \(a(\delta)\ge a_*\) is equivalent to \(\delta<(n-a_*+1)/n\).
\end{proof}

\begin{theorem}[corridor from certificates]\label[theorem]{thm:corridor}
Suppose \(L(a)\le N(a)\le U(a)\) are certified lower and upper bounds on the same staircase.  Then any first safe agreement \(a_*\) satisfies
\[
        1+
        \max\{a:L(a)>B_*\}
        \le a_*\le
        \min\{a:U(a)\le B_*\},
\]
with the evident empty-set conventions.  In particular, if for some \(a_0\)
\[
        L(a_0)>B_*\ge U(a_0+1),
\]
then \(a_*=a_0+1\).
\end{theorem}

\begin{proof}
If \(L(a)>B_*\), then \(N(a)>B_*\), so \(a\) is unsafe and the first safe point is larger.  If \(U(a)\le B_*\), then \(a\) is safe.  The one-step statement is the adjacent case.
\end{proof}

\section{The conditional frontier theorem}\label{sec:conditional}

The exact unsafe side is now settled by \cref{thm:exact-unsafe}.  The remaining safe side is isolated into three inputs.  The first controls prefix fibers.  The second controls the interior split-pencil census left by the aperiodic reduction.  The third controls primitive shift-pair terms in the second-moment ledger.  This is exactly the active-interface convention of the main raw-working file: the older broad band statement and the first-form rank-one/split-pencil assertions are retained as reductions and diagnostics, while the active closure consumes only Q, BC, and SP.

\subsection{Input Q: prefix and quotient flatness}

For \(m=K+w\), define the identity prefix map
\[
        \Phi_w:\binom Dm\longrightarrow \BB^w,
        \qquad
        M\longmapsto\bigl((-1)^j e_j(M)\bigr)_{1\le j\le w}.
\]
The lower floor uses only the average-fiber bound
\[
        \max_z|\Phi_w^{-1}(z)|\ge \frac{\binom nm}{|\BB|^w}.
\]
The safe side needs the matching upper bound after quotient-periodic structure has been absorbed at its own scale.

\begin{residualinput}[prefix/quotient flatness, Q]\label[residualinput]{prob:Q}
In the dense frontier window, prove
\[
        \max_z|\Phi_w^{-1}(z)|
        \le \kappa\binom nm|\BB|^{-w},
\]
with \(\kappa=n^{O(1)}\) in the asymptotic reserve regime and with an explicit row-dependent constant in the finite adjacent regime.  The finite constants must fit inside the margins of \cref{tab:cert-margins}.  The same statement is required at every nontrivial quotient scale after its composite directions are charged to coarser scales.
\end{residualinput}

\subsection{Input BC: base-field-normalized split-pencil census}

The non-prefix residuals can be organized by a split-pencil normal form.  In compressed language, one studies pairs of divisor-like objects whose leading windows are forced by a received line and whose remaining coefficients form a determinantal incidence problem over the base field.  The boundary profile of this census is exactly \cref{prob:Q}; the interior profiles are the genuine aperiodic band.

\begin{residualinput}[base-field-normalized split-pencil census, BC]\label[residualinput]{prob:BC}
For every interior balanced profile in the split-pencil reduction, prove that the number of primitive base-field solutions is at most its random-model count times \(e^{o(n)}\) in the asymptotic regime, and at most an explicit integer \(U_{\rm BC}(a)\) in the finite adjacent regime.  The model must include the unavoidable \(|\BB|\)-scale floor terms; quotient-pullback and common-GCD strata are removed before the primitive count is applied.
\end{residualinput}

\subsection{Input SP: primitive shift-pair control}

The exact second-moment stratification of prefix collisions contains shift-pair strata.  Quotient-pullback shift pairs are paid elsewhere.  What remains is the primitive part.

\begin{residualinput}[primitive shift-pair control, SP]\label[residualinput]{prob:SP}
For the shift-pair strata \(\mathrm{sp}_w(e;D')\) in the second-moment ledger, bound the primitive part after all quotient-pullback pairs have been removed.  The asymptotic form may lose a polynomial factor with agreement reserve.  The finite form requires explicit constants compatible with the adjacent-pair margins.
\end{residualinput}

\begin{theorem}[conditional frontier closure]\label[theorem]{thm:conditional-frontier}
Assume the asymptotic forms of \cref{prob:Q,prob:BC,prob:SP} with polynomial or \(e^{o(n)}\) losses.  Then \cref{conj:asymptotic-frontier} holds with any agreement reserve \(R\gg\log n\).  If the finite forms hold with constants fitting below the row margins in \cref{tab:cert-margins}, then the finite adjacent pairs in \cref{conj:finite-frontier} follow.
\end{theorem}

\begin{proof}
For the lower side, use the exact unsafe certificates of \cref{thm:exact-unsafe}.  For the upper side, decompose the numerator staircase by support type and field of definition.  Quotient and prefix cells are charged to \cref{prob:Q}; interior aperiodic split-pencil cells are charged to \cref{prob:BC}; primitive shift-pair cells are charged to \cref{prob:SP}; tangent cells in the high-agreement range are bounded by \cref{thm:deep-mca,prop:tangent-cell}; and below half distance the support-mismatch layer is removed by \cref{thm:mca-from-ca}.  Thus the three inputs supply an upper staircase \(U(a)\) for every agreement beyond the entropy-subfield frontier.

In the asymptotic setting, moving \(R\gg\log n\) agreements beyond the entropy crossing gives an entropy deficit \(\Omega(R)\) bits, so every polynomial loss is absorbed.  Therefore \(U(a)\le B_*\) past the frontier, while the identity-prefix floor gives \(L(a)>B_*\) before it.  \Cref{thm:corridor} localizes the threshold.

In the finite deployed setting there is no reserve.  The exact lower inequalities give \(L(a_0)>B_*\).  The finite versions of Q, BC, and SP, together with the paid tangent/quotient cells, give an explicit integer upper bound \(U(a_0+1)\).  If this fits under \(B_*\), \cref{thm:corridor} gives \(a_*=a_0+1\) exactly.  The row margins in \cref{tab:cert-margins} are precisely the allowed error budgets for these finite upper constants.
\end{proof}



\section{Safe-side decomposition into paid cells}\label{sec:cell-decomp}

The conditional theorem in \cref{sec:conditional} is useful only if the three inputs are recognizably the whole residue.  This section records the decomposition used in the proof.  It is a finite partition of possible bad-slope witnesses into cells that are already paid, and one primitive residue that becomes the split-pencil problem.

Fix an agreement integer \(m\), a code dimension \(K\in\{k,k+1\}\), and write \(w=m-K\).  A bad event comes with a support \(T\subseteq D\), \(|T|=m\), and a polynomial explanation on \(T\).  The proof ledger classifies \(T\) and the line producing it.  The classification is deliberately coarse; it is meant to be stable under finite exact verification.

\begin{definition}[paid and primitive cells]\label[definition]{def:paid-cells}
A support-line witness belongs to a paid cell if it falls into at least one of the following classes.
\begin{enumerate}[label=\textup{(\roman*)}]
\item \emph{Tangent/deep cell:} the received pair is column-close to a codeword pair, so the bad slopes are among explicit tangent ratios.  This is paid by \cref{thm:deep-mca,thm:mca-from-ca}.
\item \emph{Common-support cell:} both components interpolate on the same support.  This is not MCA-bad by definition, and in CA accounting it is a close-pair cell.
\item \emph{Quotient-pullback cell:} the support is a union of fibers of a nontrivial complete scale map.  This is charged to the quotient ledger and recursively to Q at the quotient scale.
\item \emph{Prefix-boundary cell:} the witness is a split-locator boundary profile, equivalently a fixed leading-prefix fiber of monic locators.  This is charged directly to Q.
\item \emph{Common-GCD or planted-core cell:} the split-pencil representation has a forced common divisor or planted quotient core.  This is removed before the primitive split-pencil census and charged to the appropriate lower-dimensional quotient/prefix ledger.
\item \emph{Primitive shift-pair cell:} the witness appears only in the exact second-moment expansion through a shift-pair stratum not accounted for by quotient pullback.  This is charged to SP.
\end{enumerate}
The remaining witnesses are called primitive aperiodic witnesses.
\end{definition}

\begin{proposition}[cell reduction]\label[proposition]{prop:cell-reduction}
After removing the paid cells of \cref{def:paid-cells}, every primitive aperiodic witness at agreement \(m\) is represented by a rank-one support of the received line.  The set of such supports is bounded by the base-field-normalized split-pencil census BC.
\end{proposition}

\begin{proof}
The interpolation test of \cref{lem:interp-functionals} turns support-wise explanation of \(u+zv\) into the linear equations
\[
        \alpha(T)+z\beta(T)=0.
\]
If \(\alpha(T)=\beta(T)=0\), the support is common and is not an MCA witness.  If \(T\) is quotient-periodic, it has already been charged to Q at a smaller scale.  Otherwise \(\alpha(T),\beta(T)\) are proportional and nonzero; this is the rank-one condition of \cref{thm:slope-elim-core}.  The lattice reduction \cref{prop:lattice-locus-core,thm:dichotomy-core} then converts the rank-one count into either a low-degree tangent/near-rational cell or a balanced two-generator pencil.  Low-degree cells are paid by the tangent and common-support clauses.  Balanced cells are exactly the input BC after common-GCD, planted-core, and quotient-pullback components are deleted.  The shift-pair terms do not arise from a single line census but from the second-moment expansion of Q, which is why they are isolated as SP rather than absorbed into BC.
\end{proof}

\begin{remark}[finite versus asymptotic bookkeeping]\label[remark]{rem:finite-asymp-cell}
For the asymptotic theorem the cell reduction may lose polynomial factors: an agreement reserve \(R\gg\log n\) beyond the entropy crossing absorbs them.  For the deployed adjacent pairs, no such reserve is available.  The finite forms of Q, BC, and SP must therefore return actual integer upper bounds cell by cell, and the sum of all paid and primitive cells must stay below the budget \(B_*\).  This is why \cref{conj:finite-frontier} is phrased as an adjacent-pair conjecture rather than merely as an asymptotic formula.
\end{remark}

\section{Rank-one and split-pencil reduction}\label{sec:splitpencil}

The input BC in \cref{prob:BC} is not a black box without structure.  It is the endpoint of an elementary reduction from bad slopes to a rank-one support census, then to a two-generator polynomial lattice, and finally to a split-pencil problem.  This section records that reduction in the form needed by the conditional proof.

\subsection{Interpolation functionals and slope elimination}

Let \(T\subseteq D\) have size \(m\), and put \(w=m-K\), where \(K\) is the Reed--Solomon dimension used for the list or point code.  For a word \(f:D\to\F\), define
\[
        S_r(f,T)=\sum_{x\in T}\frac{f(x)x^r}{\Lambda_T'(x)},
        \qquad 0\le r\le w-1,
\]
where \(\Lambda_T(X)=\prod_{x\in T}(X-x)\).

\begin{lemma}[interpolation functionals]\label[lemma]{lem:interp-functionals}
The restriction \(f|_T\) extends to a polynomial of degree \(<K\) if and only if
\[
        S_0(f,T)=S_1(f,T)=\cdots=S_{w-1}(f,T)=0.
\]
The functionals are \(\F\)-linear in \(f\), and are \(\BB\)-valued when \(f\) and \(T\) are \(\BB\)-valued.
\end{lemma}

\begin{proof}
The interpolant on \(T\) is
\[
        I_f(X)=\sum_{x\in T} f(x)\frac{\Lambda_T(X)}{(X-x)\Lambda_T'(x)}.
\]
It has degree at most \(m-1\).  It has degree \(<K=m-w\) exactly when its top \(w\) coefficients vanish.  The coefficient of \(X^{m-1-j}\) is a unitriangular \(T\)-symmetric linear combination of \(S_0(f,T),\ldots,S_j(f,T)\).  Hence vanishing of the top \(w\) coefficients is equivalent to vanishing of all \(S_r\), \(r<w\).  Linearity and the field-of-definition claim are immediate from the displayed formula.
\end{proof}

For a received line \((u,v)\), set
\[
        \alpha(T)=(S_r(u,T))_{r<w},
        \qquad
        \beta(T)=(S_r(v,T))_{r<w}.
\]
Call \(T\) \emph{rank-one} for \((u,v)\) if \(\alpha(T)\) and \(\beta(T)\) are proportional and not both zero, with \(\beta(T)\ne0\).  Such a support carries the unique slope \(z(T)\) satisfying \(\alpha(T)=-z(T)\beta(T)\).

\begin{theorem}[slope elimination]\label[theorem]{thm:slope-elim-core}
Fix \((u,v)\), an agreement \(m>n/2\), and \(w=m-K\).  Then:
\begin{enumerate}[label=\textup{(\alph*)}]
\item a slope \(z\) is explained on a support \(T\) of size \(m\) if and only if \(\alpha(T)+z\beta(T)=0\);
\item if some support has \(\alpha(T)=\beta(T)=0\), then \((u,v)\) itself has correlated agreement on \(T\), so no slope is MCA-bad at radius \(1-m/n\);
\item otherwise the number of MCA-bad finite slopes is at most the number of rank-one supports.
\end{enumerate}
\end{theorem}

\begin{proof}
The first item is \cref{lem:interp-functionals} applied to \(u+zv\).  The second says that if both \(u\) and \(v\) interpolate on the same support, then every linear combination does so and the pair is already mutually explained.  For the third, every bad slope must have at least one support on which it is explained; without common supports, a support carries at most one slope, and supports with \(\beta=0\ne\alpha\) carry none.
\end{proof}

Thus the aperiodic safe-side problem is a census problem for rank-one supports.  The model size for a random line is roughly \(\binom nm q^{-(w-1)}\), up to the small paid strata.  This motivates the following census target.

\begin{problem}[rank-one census]\label[problem]{prob:R1-core}
At band agreements, prove for every line \((u,v)\)
\[
        \#\mathrm{R1}(u,v;m)
        \le e^{o(n)}\max\left(1,\binom nm q^{-(w-1)}\right),
\]
with explicit finite constants in the deployed adjacent regime.  By \cref{thm:slope-elim-core}, this implies the aperiodic slope bound needed for BC, except for common-support and quotient/pullback strata, which are paid separately.
\end{problem}

\subsection{The lattice form of the census}

For a word \(U:D\to\F\), let \(\widehat U\) be its interpolating polynomial of degree at most \(n-1\), and define
\[
        M_U=\{(W,N)\in\F[X]^2: W(x)U(x)=N(x)\text{ for all }x\in D\}.
\]
Give pairs the shifted degree
\[
        \operatorname{wdeg}(W,N)=\max(\deg W,\deg N-(K-1)).
\]

\begin{proposition}[census as a lattice locus]\label[proposition]{prop:lattice-locus-core}
The module \(M_U\) is free of rank two with basis \((1,\widehat U),(0,\Lambda_D)\).  A shifted weak-Popov basis has row degrees \(d_1\le d_2\) satisfying
\[
        d_1+d_2=n-K+1.
\]
Moreover the number of \(m\)-supports on which \(U\) agrees with a degree \(<K\) polynomial is exactly
\[
        \#\{(W,N)\in M_U:
        W\text{ monic},\deg W=n-m,
        W\mid\Lambda_D,
        W\mid N\}.
\]
Equivalently this is the \(m\)-th binomial moment \(\sum_c \binom{\operatorname{agr}(U,c)}m\) of the agreement profile of \(U\) with codewords.
\end{proposition}

\begin{proof}
Every \((W,N)\in M_U\) satisfies \(N\equiv W\widehat U\pmod{\Lambda_D}\), so \((W,N)=W(1,\widehat U)-C(0,\Lambda_D)\).  The determinant of the basis is \(\Lambda_D\), of shifted degree \(n-K+1\), which gives the sum of row degrees in weak-Popov form.  If \(U\) agrees with \(c\) on \(T\), then with \(W=\Lambda_{D\setminus T}\) and \(N=Wc\) all displayed conditions hold.  Conversely such \((W,N)\) gives \(c=N/W\) of degree \(<K\), and \(W(x)(U(x)-c(x))=0\) forces agreement off the roots of \(W\).  Counting pairs \((T,c)\) gives the binomial-moment formula.
\end{proof}

\begin{theorem}[near-rational dichotomy]\label[theorem]{thm:dichotomy-core}
Let \(d_1=d_1(U)\) be the minimal shifted degree of a nonzero element of \(M_U\), and put \(\omega=n-m\).  Assume \(2w\le n-K\).
\begin{enumerate}[label=\textup{(\roman*)}]
\item If \(d_1\le w\), then either \(U\) lies within Hamming distance \(d_1\) of a unique codeword and the census is exactly \(\binom{n-d_1}{m}\), or the census is zero.
\item If \(d_1\ge w+1\), then every census element lies in a two-parameter family
\[
        Ag_1+Bg_2,
        \qquad
        \deg A\le\omega-d_1,
        \quad
        \deg B\le\omega-d_2,
\]
where \(g_1,g_2\) is the shifted basis.  This family has dimension \(\omega-w+1\).
\end{enumerate}
\end{theorem}

\begin{proof}
If \(d_1\le w\), then \(d_2=n-K+1-d_1\ge\omega+1\), so predictable degrees force every census element of shifted degree \(\omega\) to be a multiple of the first basis vector.  This gives either an error-locator pair for a unique nearby codeword or no split divisible element.  The exact count is then the choice of an \(m\)-subset inside the agreement set.  If \(d_1\ge w+1\), both basis rows may participate, and predictable degrees give precisely the displayed degree caps; the dimension is \((\omega-d_1+1)+(\omega-d_2+1)=\omega-w+1\).
\end{proof}

\begin{corollary}[reduction to the balanced core]\label[corollary]{cor:balanced-core}
For a line \((u,v)\), if two finite slopes \(z_1,z_2\) have \(d_1(u+z_iv)\le w\) with nonzero census, then the pair \((u,v)\) already has correlated agreement on a support larger than the band threshold, and no slope is MCA-bad there.  Hence, up to at most one exceptional slope, the safe-side problem lies in the balanced case \(d_1\ge w+1\).
\end{corollary}

\begin{proof}
Each low-\(d_1\) slope gives a unique nearby codeword with error of weight at most \(w\).  Subtracting two such relations shows that \(v\) is within \(2w\) of a codeword and then \(u\) is within \(3w\) of a codeword.  In the deployed band this common agreement is far larger than \(m\), so the pair is mutually explained.  If there is only one such slope, it contributes at most one to the numerator and is harmless.
\end{proof}

\subsection{Final split-pencil form}

The balanced core can be stated without the word \(U\).  A basis \((W_1,N_1),(W_2,N_2)\) of some \(M_U\) satisfies
\[
        W_1N_2-W_2N_1=\gamma\Lambda_D,
        \qquad \gamma\in\F^\times,
        \qquad \gcd(W_1,W_2)=1.
\]
Conversely, such a determinantal representation recovers \(U\) by \(U(x)=N_i(x)/W_i(x)\) for either \(i\) with \(W_i(x)\ne0\).

\begin{problem}[split-pencil problem, final form of BC]\label[problem]{prob:split-core}
For every balanced determinantal representation of \(\Lambda_D\), prove
\[
\#\left\{(A,B):
\begin{array}{l}
\deg A\le\omega-d_1,
\quad \deg B\le\omega-d_2,\\
AW_1+BW_2\text{ is monic of degree }\omega,\\
AW_1+BW_2\mid\Lambda_D
\end{array}\right\}
\le e^{o(n)}\max\left(1,\binom n\omega q^{1-w}\right),
\]
with explicit finite constants in the deployed adjacent regime.  For multiplicative rows \(\Lambda_D=X^n-\beta\).  For circle rows \(\Lambda_D\) is the twin-coset domain polynomial.  This problem implies \cref{prob:R1-core}, hence BC.
\end{problem}

The importance of \cref{prob:split-core} is that the whole aperiodic residue has become a single algebraic counting problem: divisibility by the domain polynomial inside a two-generator pencil.  Common-GCD, quotient-pullback, and sunflower/planted strata are visible inside the coordinates \((A,B)\) and are paid before the primitive count is applied.


\subsection{Base-field normalization and subfield floors}\label{subsec:base-field-normalization}

The census model in \cref{prob:split-core} must be normalized at the base-field scale whenever the evaluation domain lies in a proper subfield.  This point is not cosmetic.  The identity-prefix witnesses themselves produce large base-field split-pencil cells over extension fields, so a purely \(q\)-random model would be false in the deployed KoalaBear row.

Let \(p=|\BB|\), \(q=|\F|\), and assume \(D\subseteq\BB\).  At agreement \(m\), with \(w=m-K\), the boundary profile of the split pencil has size about
\[
        \binom nm p^{-w},
\]
not \(\binom nm q^{-w}\).  Interior profiles inherit the same phenomenon.  If \(d_1\ge w+1\) is the first shifted degree and \(m'=K-1+d_1\), the natural base-field floor is
\begin{equation}\label{eq:MB-floor}
        M_{\BB}(d_1;m)
        :=
        \binom{m'}m\binom n{m'}p^{-(d_1-1)},
        \qquad m'=K-1+d_1.
\end{equation}
The factor \(\binom{m'}m\) records the choice of an \(m\)-support inside the larger level-\(m'\) prefix witness.

\begin{proposition}[subfield floors for split-pencil censuses]\label[proposition]{prop:subfield-census-floor}
Assume \(D\subseteq\BB\), \(|D|=n\), and \(K\le m\).  Let \(w=m-K\).  Then:
\begin{enumerate}[label=\textup{(\alph*)}]
\item At the boundary profile \(d_1=w+1\), the heaviest identity-prefix witness has census at least
\[
        \left\lceil \binom nm p^{-w}\right\rceil.
\]
\item For every interior profile \(d_1\ge w+2\) satisfying \(m'=K-1+d_1\le n\), there is a base-field witness whose level-\(m\) census is at least \(M_{\BB}(d_1;m)\), up to the harmless integer ceiling in the second binomial factor.
\item For the non-base pole lines used in the flexible conversion, every support in the heaviest boundary fiber is a non-common rank-one support.  Thus the per-line \(\Rone\) census also has a base-field floor of order \(\binom nm p^{-w}\).
\end{enumerate}
\end{proposition}

\begin{proof}
The first item is exactly \cref{lem:identity-floor,prop:exact-witness-list}.  For the second, run the same prefix construction at level \(m'\), fixing \(d_1-1=m'-K\) leading coefficients.  The heaviest level-\(m'\) fiber has at least \(\binom n{m'}p^{-(d_1-1)}\) members.  Each listed codeword agrees on exactly its \(m'\)-root set, and therefore contributes \(\binom{m'}m\) size-\(m\) supports to the binomial moment census; distinct level-\(m'\) locators give distinct codewords, so the supports are counted without unintended merging.  For the third item, take a pole \(\alpha\in\F\setminus\BB\).  If \(T\) is a prefix-fiber support with locator \(\Lambda_T\), the associated close slope is obtained from the value \(\Lambda_T(\alpha)\).  The second component \(-1/(x-\alpha)\) cannot interpolate on \(T\) by a degree \(<K\) polynomial when \(|T|\ge K+1\), because \((X-\alpha)G(X)+1\) would have too many roots and value \(1\) at \(\alpha\).  Hence the support is non-common and rank-one by \cref{thm:slope-elim-core}.
\end{proof}

\begin{problem}[base-field-normalized split-pencil census, expanded BC]\label[problem]{prob:BC-expanded}
The finite and asymptotic forms of BC should bound primitive balanced split-pencil cells by the maximum of the challenge-field random term and the base-field floors that are actually present:
\[
        e^{o(n)}
        \max\left(
             1,
             M_{\BB}(d_1;m),
             \binom n\omega q^{1-w}
        \right)
\]
for asymptotic profiles, with \(e^{o(n)}\) replaced by explicit constants in the finite deployed adjacent regime.  The primitive qualifier means that quotient-pullback, common-GCD, planted-core, and boundary-prefix cells are removed before the bound is applied.
\end{problem}

\begin{remark}[why the normalization matters]
In an extension row such as KoalaBear, \(q=p^6\).  A model with only \(q\)-scale randomness would undercount the base-field prefix cells by a factor exponential in \(w\).  The corrected BC input does not weaken the program; it states the model at the scale where true witnesses live and asks for randomness only after the unavoidable subfield floors have been paid.
\end{remark}

\subsection{Primitive shift-pair control}\label{subsec:sp-expanded}

The input SP is the smallest of the three named inputs conceptually, but it is crucial in finite margins.  It appears when the second moment of prefix fibers is decomposed more finely than the balanced-ternary ledger of \cref{prop:moment-kernel}.  Some collisions arise from shifting or pulling back a smaller quotient configuration; those have already been charged.  The primitive residue is the one not explained by any quotient symmetry.

\begin{definition}[primitive shift pair]\label[definition]{def:primitive-shift-pair}
A shift-pair collision at depth \(w\) is a pair of disjoint sets \((A,A')\) with equal first \(w\) moments,
\[
        \sum_{x\in A}x^j=\sum_{x\in A'}x^j,
        \qquad 1\le j\le w,
\]
together with a nontrivial affine or multiplicative relation carrying one part of the support to the other inside a quotient scale.  It is primitive if no proper quotient map, common divisor, or complete fiber pullback explains the relation.
\end{definition}

\begin{problem}[SP, expanded form]\label[problem]{prob:SP-expanded}
For every depth \(w\), agreement \(m\), and quotient scale appearing in the prefix ledger, prove that the primitive shift-pair contribution to the relevant collision moment is bounded by the random-model contribution times a polynomial loss asymptotically, and by an explicit integer constant in the finite adjacent rows.
\end{problem}

\begin{remark}[SP as an audit item]
SP is included separately because it is easy to lose it in prose.  The prefix moment ledger can appear closed after Q if one only controls maximum fibers, but the finite proof needs the exact decomposition of second and higher moments.  Primitive shift pairs are the residual configurations whose quotient ancestor is not visible in the first-order prefix count.
\end{remark}


\section{Dual-weighted syndrome calculus}\label{sec:syndrome-calculus}

The slope-elimination proof in \cref{sec:splitpencil} used the unweighted interpolation functionals \(S_r(f,T)\).  For exact finite audits it is often better to use the dual-weighted normalization attached to the full domain polynomial.  This section records the equivalent form.  It is useful for two reasons.  First, the moment windows become ordinary syndrome windows.  Second, the split-pencil equations become rank conditions on Hankel matrices, which is the form in which finite certificates can be audited.

Let
\[
        P_D(X)=\prod_{x\in D}(X-x),
        \qquad
        v_x=\frac{1}{P_D'(x)}.
\]
The weights \(v_x\) are nonzero because the evaluation points are distinct.

\begin{lemma}[residue pairing]\label[lemma]{lem:residue-pairing}
For every polynomial \(h\in\F[X]\) of degree at most \(n-1\),
\[
        \sum_{x\in D}v_xh(x)=[X^{n-1}]h(X).
\]
In particular the sum is zero whenever \(\deg h\le n-2\).
\end{lemma}

\begin{proof}
The Lagrange formula gives
\[
        h(X)=\sum_{x\in D} h(x)\frac{P_D(X)}{(X-x)P_D'(x)}
\]
for \(\deg h\le n-1\).  The polynomial \(P_D(X)/(X-x)\) is monic of degree \(n-1\), so comparing the coefficient of \(X^{n-1}\) gives the displayed identity.
\end{proof}

\begin{definition}[weighted syndrome window]\label[definition]{def:weighted-syndrome}
For a word \(y:D\to\F\), define
\[
        s_j(y)=\sum_{x\in D}v_xy(x)x^j,
        \qquad j\ge0.
\]
For a polynomial \(L(X)=\sum_i\ell_iX^i\), say that \(L\) annihilates \(y\) in a window of length \(\tau\), written \(L\in\mathcal A_\tau(y)\), if
\[
        \sum_i\ell_i s_{i+t}(y)=0,
        \qquad 0\le t\le\tau-1.
\]
Equivalently,
\[
        \sum_{x\in D}v_xy(x)L(x)x^t=0
        \qquad(0\le t<\tau).
\]
\end{definition}

\begin{lemma}[syndrome interpolation criterion]\label[lemma]{lem:syndrome-criterion}
Let \(T\subseteq D\) have size \(m\), and let \(E=D\setminus T\) have locator \(L_E(X)=\prod_{x\in E}(X-x)\), of degree \(\omega=n-m\).  A word \(y:D\to\F\) agrees on \(T\) with a polynomial of degree \(<K\) if and only if
\[
        L_E\in\mathcal A_{m-K}(y).
\]
\end{lemma}

\begin{proof}
Suppose \(y=c\) on \(T\), with \(\deg c<K\).  The word \(L_E(y-c)\) vanishes on all of \(D\).  Hence for \(0\le t<m-K\),
\[
        \sum_{x\in D}v_x y(x)L_E(x)x^t
        =\sum_{x\in D}v_x c(x)L_E(x)x^t.
\]
The degree of \(cL_EX^t\) is at most \((K-1)+\omega+(m-K-1)=n-2\), so the right side is zero by \cref{lem:residue-pairing}.  Thus \(L_E\) annihilates \(y\) in the claimed window.

Conversely, assume the annihilation equations.  Interpolate \(y|_T\) by a polynomial \(I\) of degree \(<m\).  Since \(L_E(y-I)\) vanishes on \(D\), the same residue computation shows that the top \(m-K\) coefficients of \(I\) vanish.  Hence \(\deg I<K\), so \(I\) is the desired polynomial explanation.
\end{proof}

For a line \((u,v)\), the condition that \(u+zv\) be explained on \(T\) becomes
\[
        L_E\in\mathcal A_w(u+zv),
        \qquad w=m-K.
\]
Equivalently the syndrome vectors
\[
        \sigma_u(L_E)=(\sum_i\ell_i s_{i+t}(u))_{0\le t<w},
        \qquad
        \sigma_v(L_E)=(\sum_i\ell_i s_{i+t}(v))_{0\le t<w}
\]
are proportional.  This is the same rank-one condition as in \cref{thm:slope-elim-core}, conjugated by the diagonal dual weights and a triangular change of basis.

\begin{proposition}[Hankel rank-one form]\label[proposition]{prop:hankel-rank-one}
Fix \(\omega=n-m\) and \(w=m-K\).  Write \(L_E(X)=X^\omega+\ell_1X^{\omega-1}+\cdots+\ell_\omega\).  For a word \(y\), the annihilation equations \(L_E\in\mathcal A_w(y)\) are affine linear equations in \((\ell_1,\ldots,\ell_\omega)\) whose coefficient matrix is the \(w\times\omega\) Hankel window
\[
        H_{w,\omega}(y)=\bigl(s_{\omega-i+t}(y)\bigr)_{0\le t<w,\,1\le i\le\omega}.
\]
For a line \((u,v)\), support explanations at agreement \(m\) are contained in the rank-one locus of the two affine windows
\[
        b_u+H_{w,\omega}(u)\ell,
        \qquad
        b_v+H_{w,\omega}(v)\ell,
\]
where \(b_y=(s_{\omega+t}(y))_{0\le t<w}\).
\end{proposition}

\begin{proof}
Expanding \(L_E\) in the equations of \cref{def:weighted-syndrome} gives
\[
        s_{\omega+t}(y)+\sum_{i=1}^\omega\ell_i s_{\omega-i+t}(y)=0,
        \qquad 0\le t<w.
\]
This is precisely the displayed affine Hankel system.  For a line, the system for \(u+zv\) is the linear combination of the two affine systems, so existence of \(z\) is exactly proportionality of the two syndrome-window vectors, unless both vanish, which is the common-support cell.
\end{proof}

\begin{remark}[why the dual weights matter]
Over a multiplicative coset one can often hide the dual weights in a monomial shift, but that convention introduces a constant-term defect at the boundary of the window.  The residue normalization has no defect and works for arbitrary \(D\), including circle and map-smooth rows.  Since multiplying coordinates by nonzero weights is a diagonal equivalence of generalized Reed--Solomon rows, Hamming supports and MCA numerators are unchanged.
\end{remark}

\subsection{From Hankel windows to split pencils}

The lattice \(M_U\) of \cref{prop:lattice-locus-core} can be recovered from the syndrome windows.  If \(U\) is a word and \(L\) is a degree-\(\omega\) locator, then \(L\in\mathcal A_w(U)\) exactly says that there exists a degree \(<K\) numerator \(N\) with \(LU=N\) on the complement of the roots of \(L\).  Thus the solution variety of the Hankel equations is the projection of the lattice locus in \cref{prop:lattice-locus-core}.  In balanced degree, a weak-Popov basis \((W_1,N_1),(W_2,N_2)\) gives every candidate denominator as
\[
        L=AW_1+BW_2.
\]
The split-pencil census asks how often such an \(L\) divides \(P_D\).  In the multiplicative case \(P_D=X^n-\beta\), so this is a divisor problem inside a two-dimensional affine polynomial pencil.  In a circle row, \(P_D\) is the twin-coset coordinate polynomial, and the same formulation remains valid.

\begin{proposition}[syndrome-to-pencil equivalence]\label[proposition]{prop:syndrome-pencil-equivalence}
For fixed \(U\), agreement \(m\), and \(w=m-K\), the following counts are equal:
\begin{enumerate}[label=\textup{(\alph*)}]
\item pairs \((T,c)\) with \(|T|=m\), \(c\in\RS[\F,D,K]\), and \(U=c\) on \(T\);
\item monic divisors \(L\mid P_D\), \(\deg L=n-m\), with \(L\in\mathcal A_w(U)\), counted with the numerator determined by interpolation;
\item lattice elements \((W,N)\in M_U\) satisfying the divisibility conditions of \cref{prop:lattice-locus-core}.
\end{enumerate}
\end{proposition}

\begin{proof}
Given \((T,c)\), let \(L=L_{D\setminus T}\).  Then \cref{lem:syndrome-criterion} gives \(L\in\mathcal A_w(U)\), and \((L,Lc)\in M_U\).  Conversely a monic divisor \(L\mid P_D\) determines its zero set \(E\subseteq D\), hence \(T=D\setminus E\); the annihilation condition gives a degree \(<K\) interpolant on \(T\).  The equivalence with lattice elements is exactly the identity \(N\equiv W\widehat U\pmod{P_D}\), with \(W=L\).
\end{proof}

This calculus gives the finite certificate interface: a computer certificate may list Hankel-rank charts, weak-Popov profiles, and divisor tests in the pencil, while the paper reasons invariantly in terms of supports and lattices.


\section{Why Q is the central measurable problem}\label{sec:Q}

This section records the self-contained dictionary behind \cref{prob:Q}.  It explains both why prefix flatness is plausible and why low moments alone do not settle the finite prize.

Assume for simplicity that \(0\notin D\).  For \(w\ge1\), define the moment kernel
\[
        K_w=\left\{f:D\to\BB:\sum_{x\in D}f(x)x^j=0\ \text{ for }1\le j\le w\right\}.
\]
When \(w<\operatorname{char}\BB\), Newton identities give a triangular bijection between the first \(w\) elementary symmetric functions and the first \(w\) power sums.  Thus equal prefix fibers can be studied through power-sum equalities.

\begin{proposition}[moment-kernel dictionary]\label[proposition]{prop:moment-kernel}
The code \(K_w\) has dimension \(n-w\) and minimum Hamming weight at least \(w+1\).  If two distinct \(m\)-subsets \(M,M'\) lie in the same prefix fiber, then
\[
        |M\triangle M'|\ge w+1.
\]
Moreover, if \(T_w(\ell)\) denotes the number of ordered pairs of disjoint \(\ell\)-subsets \(A,A'\subseteq D\) with equal first \(w\) power sums, then
\[
        \sum_z|\Phi_w^{-1}(z)|^2
        =\binom nm+\sum_{\ell\ge\ceil{(w+1)/2}}\binom{n-2\ell}{m-
\ell}T_w(\ell).
\]
Thus the second moment of prefix fibers is a weighted balanced-ternary weight enumerator of \(K_w\).
\end{proposition}

\begin{proof}
The map \(f\mapsto(\sum_x f(x)x^j)_{1\le j\le w}\) has full rank \(w\), since any \(s\le w\) columns form a Vandermonde minor with determinant \((\prod_i x_i)\prod_{i<j}(x_j-x_i)\ne0\).  This proves the dimension and minimum-weight statement.  If \(M,M'\) have the same first \(w\) power sums, then \(\one_M-\one_{M'}\in K_w\), so its nonzero support, \(M\triangle M'\), has size at least \(w+1\).

For the second moment, count ordered pairs \((M,M')\) in the same fiber.  Write \(A=M\setminus M'\), \(A'=M'\setminus M\), with \(|A|=|A'|=\ell\).  The common part has size \(m-\ell\) and may be any subset of the remaining \(n-2\ell\) points.  Equality of prefixes is exactly equality of the first \(w\) power sums of \(A\) and \(A'\).  The term \(\ell=0\) gives \(\binom nm\); all other terms have \(2\ell=|M\triangle M'|\ge w+1\).
\end{proof}

\begin{corollary}[first all-depth fiber bound]\label[corollary]{cor:packing-bound}
For every prefix value \(z\),
\[
        |\Phi_w^{-1}(z)|
        \le
        \frac{\binom n{m-
        \delta+1}}{\binom m{\delta-1}},
        \qquad
        \delta=\ceil{(w+1)/2}.
\]
\end{corollary}

\begin{proof}
Distinct members of a fiber share at most \(m-
\delta\) points.  Hence no \((m-
\delta+1)\)-subset of \(D\) is contained in two fiber members.  Each member contains \(\binom m{\delta-1}\) such subsets, while \(D\) contains \(\binom n{m-
\delta+1}\) of them.
\end{proof}

This bound is nontrivial at full frontier depth but still far above the conjectured average.  It shows the right obstruction: a counterexample to Q must create a very large constant-weight packing inside a high-distance moment kernel.

\subsection{Collision gaps}

Let \(\mu\) be the distribution of \(\Phi_w(M)\) for a uniform \(m\)-subset \(M\subseteq D\).  Define
\[
        \Gamma_r(m,w)=|\BB|^{w(r-1)}\sum_z\mu(z)^r,
        \qquad r\ge2.
\]
The uniform distribution has \(\Gamma_r=1\).  Let
\[
        R=|\BB|^w\max_z\mu(z)
\]
be the max-to-mean fiber ratio.

\begin{proposition}[moment sandwich]\label[proposition]{prop:moment-sandwich}
For every \(r\ge2\),
\[
        \frac{\log_2\Gamma_r}{r-1}
        \le \log_2 R
        \le \frac{\log_2\Gamma_r+w\log_2|\BB|}{r}.
\]
Consequently \(R\le e^{o(n)}\) is equivalent, asymptotically, to \(\Gamma_r=e^{o(n)}\) for every fixed \(r\).  For a finite adjacent pair with a \(\Delta\)-bit margin, moments alone need a hierarchy up to roughly \(r\approx w\log_2|\BB|/\Delta\).
\end{proposition}

\begin{proof}
Since \(\mu(z)\le R|\BB|^{-w}\),
\[
        \Gamma_r\le |\BB|^{w(r-1)}(R|\BB|^{-w})^{r-1}\sum_z\mu(z)=R^{r-1}.
\]
For the other direction,
\[
        \max_z\mu(z)\le\left(\sum_z\mu(z)^r\right)^{1/r}
        =\left(\Gamma_r |\BB|^{-w(r-1)}\right)^{1/r},
\]
so \(R\le(\Gamma_r|\BB|^w)^{1/r}\).  Taking logs gives the inequalities.
\end{proof}

\begin{remark}[heuristic and falsification]
The working heuristic is that, after quotient-periodic directions are removed, \(\Phi_w\) behaves like a random constant-weight syndrome map on a moment curve.  This predicts polynomial max-to-mean inflation in the dense bulk.  A practical falsification test is therefore to measure \(R\) and the hierarchy \(\Gamma_r\), separating primitive frequencies from divisor-lattice frequencies.  Any counterexample to the frontier should be classifiable: either it produces a super-polynomial primitive prefix fiber, refuting Q, or it produces a super-polynomial primitive split-pencil/shift-pair family, refuting BC or SP.
\end{remark}

\subsection{Stable Fourier measurement}

For computational evidence one can measure Fourier coefficients without forming huge elementary symmetric coefficients.  If \(|\varepsilon_i|\le1\), define
\[
        q_j^{(i)}=\frac{1}{\binom ij}\sum_{1\le r_1<\cdots<r_j\le i}\varepsilon_{r_1}\cdots\varepsilon_{r_j}.
\]
Then
\[
        q_j^{(i)}=\left(1-\frac{j}{i}\right)q_j^{(i-1)}+\frac{j}{i}\varepsilon_i q_{j-1}^{(i-1)},
        \qquad |q_j^{(i)}|\le1.
\]
For an additive character \(\psi\) and a frequency vector \(\tau\), taking
\[
        \varepsilon_x=\psi\left(\sum_{r=1}^w\tau_rx^r\right)
\]
computes the normalized Fourier coefficient of the constant-weight power-sum distribution.  The recursion is non-expansive, so numerical error does not amplify.  The important caveat is that individual Fourier maxima are the wrong endpoint: Q requires cancellation across the whole joint tail, not just a per-frequency Weil bound.

\section{Prefix-collision rigidity and higher moments}\label{sec:prefix-collision-rigidity}

The input Q is a maximum-fiber statement.  The most concrete attack on it is through collision moments.  This section expands the dictionary of \cref{sec:Q} in the form needed for a proof attempt.  The key point is that equal leading elementary prefixes are much more rigid than an arbitrary random linear syndrome: subtracting two locators converts prefix equality into a high-degree vanishing statement, and the residual collision must be visible in a balanced moment kernel.

\subsection{Elementary prefixes and Newton windows}

For an \(m\)-subset \(M\subseteq D\), write
\[
        \Lambda_M(X)=X^m+c_1(M)X^{m-1}+\cdots+c_m(M),
        \qquad c_j(M)=(-1)^je_j(M).
\]
The prefix map used in the unsafe construction is
\[
        \Phi_w(M)=(c_1(M),\ldots,c_w(M)).
\]
When \(w<\operatorname{char}\BB\), Newton identities give an invertible triangular transformation between \((c_1,\ldots,c_w)\) and the power sums
\[
        p_j(M)=\sum_{x\in M}x^j,
        \qquad 1\le j\le w.
\]
Thus the moment-kernel formalism is exactly equivalent to elementary-prefix equality in the deployed rows, where the frontier depth \(w\approx 6.75\cdot 10^4\) is much smaller than the characteristic \(p\approx2^{31}\).

\begin{lemma}[locator subtraction]\label[lemma]{lem:locator-subtraction}
If \(M,M'\subseteq D\) are distinct \(m\)-subsets with \(\Phi_w(M)=\Phi_w(M')\), then
\[
        \Lambda_M(X)-\Lambda_{M'}(X)
\]
has degree at most \(m-w-1\).  In particular it vanishes on \(M\cap M'\), and if \(|M\cap M'|>m-w-1\), then \(M=M'\).
\end{lemma}

\begin{proof}
The leading monic terms and the next \(w\) coefficients cancel by hypothesis, leaving degree at most \(m-w-1\).  Every point in \(M\cap M'\) is a root of both locators, so it is a root of the difference.  If a nonzero polynomial of degree at most \(m-w-1\) has more than \(m-w-1\) roots, it is zero.  Equal monic locators have equal root sets.
\end{proof}

\begin{corollary}[side-distance sharpening]\label[corollary]{cor:side-distance}
Two distinct members of one elementary-prefix fiber satisfy
\[
        |M\setminus M'|=|M'\setminus M|\ge w+1.
\]
Consequently every prefix fiber is a constant-weight code of side-distance at least \(w+1\).
\end{corollary}

\begin{proof}
If \(|M\setminus M'|\le w\), then \(|M\cap M'|\ge m-w\), contradicting \cref{lem:locator-subtraction}.  The two one-sided differences have equal size because \(|M|=|M'|\).
\end{proof}

\begin{remark}[relation to the moment-kernel distance]
The moment-kernel argument in \cref{prop:moment-kernel} gives total symmetric-difference distance at least \(w+1\), hence side-distance at least \(\ceil{(w+1)/2}\).  Locator subtraction improves this to side-distance \(w+1\) for elementary-prefix fibers.  The improvement is important for finite packing bounds, but it does not by itself prove Q; a packing bound at this distance is still far above the average fiber in the dense frontier window.
\end{remark}

\subsection{Exact second moment}

Let \(F_z=\Phi_w^{-1}(z)\).  The second moment is
\[
        \sum_z |F_z|^2.
\]
Counting pairs by their one-sided difference size gives an exact formula.  Let \(C_w(e)\) be the number of ordered disjoint pairs \((A,A')\) of \(e\)-subsets of \(D\) for which there exists at least one common core \(R\subseteq D\setminus(A\cup A')\) such that
\[
        \Phi_w(R\cup A)=\Phi_w(R\cup A').
\]
This definition still depends on the core, because elementary symmetric functions are not additive.  In the power-sum coordinates the core disappears: the corresponding count \(T_w(e)\) is simply equality of the first \(w\) power sums of \(A\) and \(A'\).  Since the two coordinate systems are triangular for \(w<p\), the two ledgers are equivalent after the Newton transformation, and the power-sum form is usually the cleaner one.

\begin{proposition}[second moment, side-distance form]\label[proposition]{prop:second-moment-side}
Assume \(w<\operatorname{char}\BB\).  Let \(T_w(e)\) be the number of ordered pairs \((A,A')\) of disjoint \(e\)-subsets of \(D\) satisfying
\[
        \sum_{x\in A}x^j=\sum_{x\in A'}x^j,
        \qquad 1\le j\le w.
\]
Then
\[
        \sum_z |F_z|^2
        =
        \binom nm+
        \sum_{e\ge w+1}\binom{n-2e}{m-e}T_w(e),
\]
where the lower limit may be replaced by \(\ceil{(w+1)/2}\) if one uses only the BCH-style moment-kernel distance and by \(w+1\) when the elementary-prefix side-distance of \cref{cor:side-distance} is enforced.
\end{proposition}

\begin{proof}
Count ordered pairs \((M,M')\) in the same prefix fiber.  The diagonal pairs contribute \(\binom nm\).  For an off-diagonal pair let \(A=M\setminus M'\), \(A'=M'\setminus M\), and \(R=M\cap M'\).  Then \(|A|=|A'|=e\), the sets \(A,A'\) are disjoint, and \(R\) is any \((m-e)\)-subset of the remaining \(n-2e\) points.  In power-sum coordinates, equality of prefixes is exactly equality of the power sums of \(A\) and \(A'\), since the common core cancels.  Conversely, such \((A,A')\) and \(R\) produce a colliding ordered pair.  The lower bound on \(e\) is the side-distance statement.
\end{proof}

\subsection{Higher fixed moments}

For \(r\ge2\), the normalized collision moment was defined in \cref{prop:moment-sandwich} by
\[
        \Gamma_r=|\BB|^{w(r-1)}\sum_z\mu(z)^r,
\]
where \(\mu\) is the distribution of \(\Phi_w(M)\).  A fixed-order proof of Q would bound \(\Gamma_r\) by \(e^{o(n)}\) for every fixed \(r\).  The following formulation makes this a finite algebraic counting problem.

\begin{problem}[fixed-moment Q]\label[problem]{prob:fixed-moment-Q}
For every fixed \(r\ge2\), bound the number of \(r\)-tuples \((M_1,\ldots,M_r)\) of \(m\)-subsets of \(D\) with
\[
        \Phi_w(M_1)=\cdots=\Phi_w(M_r)
\]
by
\[
        e^{o(n)}\binom nm^r|\BB|^{-w(r-1)}
\]
in the dense frontier window after quotient-pullback tuples are removed.
\end{problem}

\begin{proposition}[fixed moments imply asymptotic Q]\label[proposition]{prop:fixed-moments-imply-Q}
If \cref{prob:fixed-moment-Q} holds for every fixed \(r\), then the asymptotic form of Q holds with \(e^{o(n)}\) max-to-mean loss.
\end{proposition}

\begin{proof}
The hypothesis is precisely \(\Gamma_r=e^{o(n)}\) for every fixed \(r\).  By the right side of the moment sandwich \cref{prop:moment-sandwich}, for any fixed \(r\)
\[
        \log R\le \frac{\log\Gamma_r+w\log|\BB|}{r}.
\]
Take \(r\) large after \(n\) but still fixed relative to \(n\).  Since \(w\log|\BB|=O(n)\), this gives \(\log R\le o(n)+O(n/r)\).  Letting \(r\to\infty\) along constants gives \(\log R=o(n)\).
\end{proof}

\begin{remark}[finite moment orders at the deployed rows]\label[remark]{rem:finite-moment-orders}
The same inequality explains the finite difficulty.  At the active adjacent
rows, the bit scale \(w\beta\) is about \(2.09\cdot10^6\) bits.  Hence a direct
moment-to-maximum conversion requires a finite hierarchy whose depth depends on
the remaining bit margin:
\[
\begin{array}{c|c|c|c}
\text{row} & w & \Delta\text{ bits} & w\beta/\Delta \\
\hline
\text{KoalaBear MCA} & 67470 & 22.2 & 9.4\cdot10^4 \\
\text{KoalaBear list} & 67470 & 22.0 & 9.5\cdot10^4 \\
\text{Mersenne-31 MCA} & 67446 & 3.3 & 6.3\cdot10^5 \\
\text{Mersenne-31 list} & 67446 & 3.1 & 6.7\cdot10^5 .
\end{array}
\]
Thus the finite adjacent rows require either direct constant-factor fiber
control or a finite hierarchy reaching the displayed order.  This is the main
raw-working source's correction to the informal slogan that order \(10^5\)
moments would be uniformly enough.
\end{remark}

\subsection{Primitive collision taxonomy}

The collision ledger decomposes \(T_w(e)\) into structured and primitive pieces.  The structured pieces are those induced by complete fibers of a quotient map, by a common divisor in the split-locator representation, or by an affine symmetry of the domain.  The primitive piece is what remains after all such explanations are removed.

\begin{definition}[primitive prefix collision]\label[definition]{def:primitive-prefix-collision}
A pair \((A,A')\) counted by \(T_w(e)\) is primitive if there is no nontrivial complete scale map \(\varphi:D\to Q\) and no pair \((B,B')\) in \(Q\) such that \(A=\varphi^{-1}(B)\), \(A'=\varphi^{-1}(B')\), up to deletion of a common planted fiber core; and if no nontrivial common factor or domain symmetry explains the equality of moments.
\end{definition}

\begin{problem}[primitive collision bound]\label[problem]{prob:primitive-collision-bound}
Prove that primitive prefix collisions satisfy the random-model bound, with polynomial loss asymptotically and exact constants in the deployed adjacent rows.  Together with quotient recursion, this is the collision-moment form of Q and SP.
\end{problem}

This taxonomy is meant to prevent double counting.  Quotient-pullback collisions are not failures of Q at the identity scale; they are lower-scale Q instances.  Shift-pair collisions are not failures of BC; they are SP instances.  A correct finite proof has to mark these classes explicitly before summing the numerator.


\section{Finite proof obligations for the adjacent pairs}\label{sec:finite-obligations}

The asymptotic frontier tolerates polynomial slack; the Proximity Prize adjacent pairs do not.  This section spells out the finite arithmetic shape of a complete proof.  It is the intended interface between the conceptual inputs Q, BC, SP and a machine-checkable certificate.

For each row, let \(a_0\) be the proved unsafe agreement and \(a_1=a_0+1\) the conjectured first safe agreement.  Let \(Q_{
m slope}=q\) for MCA rows and \(Q_{
m list}=q\) for list rows, and let
\[
        B_* = \floor{\varepsilon^* q}.
\]
A finite safe certificate at \(a_1\) consists of an explicit partition of all possible numerator contributions into paid cells and primitive cells,
\[
        N(a_1)\le
        N_{\rm deep}+N_{\rm tan}+N_{\rm quot}+N_Q+N_{BC}+N_{SP},
\]
with each summand an integer and with total at most \(B_*\).  The deployed margins in \cref{tab:cert-margins} are not margins for the entire theorem; they are the remaining room after the identity-prefix lower floor fails at \(a_1\).  In particular the Mersenne-31 rows leave only about three bits, so the finite proof must be close to exact.

\begin{definition}[finite Q certificate]\label[definition]{def:finite-Q-cert}
At a row \((D,\BB,\F,k,n)\) and agreement \(a\), a finite Q certificate is a list of quotient scales \(\varphi_i:D\to Q_i\), prefix depths \(w_i\), and integers \(U_{Q,i}\) such that every prefix or quotient-pullback cell at agreement \(a\) is assigned to one of the entries and
\[
        \max_z |\Phi_{i,w_i}^{-1}(z)|\le U_{Q,i}
\]
for the corresponding prefix map.  The total prefix contribution is the sum of the assigned \(U_{Q,i}\), with multiplicities stated by the scale compiler.
\end{definition}

\begin{definition}[finite BC certificate]\label[definition]{def:finite-BC-cert}
A finite BC certificate at agreement \(a\) is a list of balanced split-pencil profiles \((d_1,d_2,\omega)\), together with explicit integers \(U_{BC}(d_1,d_2,\omega)\), such that every primitive balanced aperiodic witness belongs to exactly one profile and the divisibility count in \cref{prob:split-core,prob:BC-expanded} is at most the stated integer for that profile.
\end{definition}

\begin{definition}[finite SP certificate]\label[definition]{def:finite-SP-cert}
A finite SP certificate is an explicit bound \(U_{SP}\) for the primitive shift-pair contribution after all quotient-pullback and common-divisor shift pairs have been deleted from the collision ledger.  It must specify the deletion rules; otherwise the same family could be counted both in Q and SP.
\end{definition}

\begin{theorem}[finite adjacent-pair compiler]\label[theorem]{thm:finite-adjacent-compiler}
Fix one of the four deployed rows in \cref{conj:finite-frontier}.  Suppose:
\begin{enumerate}[label=\textup{(\roman*)}]
\item the exact unsafe inequality of \cref{thm:exact-unsafe} holds at \(a_0\);
\item the finite Q, BC, and SP certificates at \(a_1=a_0+1\), together with the tangent/deep and quotient paid cells, give an integer upper bound \(U(a_1)\le B_*\);
\item the cell partition is disjoint or is accompanied by an inclusion-exclusion convention that only overcounts downward.
\end{enumerate}
Then \(a_1\) is the first safe agreement for that row.  Under the closed-ball convention,
\[
        \delta_C^*(\varepsilon^*)=\frac{n-a_0}{n},
\]
with the endpoint itself unsafe.
\end{theorem}

\begin{proof}
The first assumption gives a lower numerator \(L(a_0)>B_*\), hence \(a_0\) is unsafe.  The second and third assumptions give a true upper numerator \(N(a_1)\le U(a_1)\le B_*\), hence \(a_1\) is safe.  Monotonicity of the numerator in the agreement threshold and \cref{thm:corridor} imply that \(a_1\) is the unique first safe agreement.  The endpoint identity follows from \cref{thm:staircase}: all radii with \(\ceil{(1-\delta)n}\ge a_1\) are safe, while \(\delta=(n-a_0)/n\) has agreement threshold \(a_0\) and is unsafe.
\end{proof}

\begin{center}\small
\begin{tabular}{@{}lrrrr@{}}
\toprule
row & \(a_0\) & \(a_1\) & target budget & next-step room \\
\midrule
KoalaBear MCA & 1116047 & 1116048 & \(\floor{p^6/2^{128}}\) & \(22.2\) bits \\
KoalaBear list & 1116046 & 1116047 & \(\floor{p^6/2^{128}}\) & \(22.0\) bits \\
Mersenne-31 MCA & 1116023 & 1116024 & \(\floor{(p')^4/2^{100}}\) & \(3.3\) bits \\
Mersenne-31 list & 1116022 & 1116023 & \(\floor{(p')^4/2^{100}}\) & \(3.1\) bits \\
\bottomrule
\end{tabular}
\end{center}

\begin{remark}[why polynomial finite losses are insufficient]
At \(n=2^{21}\), even a factor of \(n\) costs \(21\) bits.  Such a loss would already exceed the entire Mersenne-31 adjacent margin.  This is why the asymptotic and finite statements are separated throughout the paper.  Polynomial-loss Q/BC/SP settles the asymptotic frontier with logarithmic reserve; exact-constant Q/BC/SP is needed for the printed adjacent pairs.
\end{remark}

\begin{problem}[finite settlement checklist]\label[problem]{prob:finite-checklist}
For each of the four rows, produce the following certificate files.
\begin{enumerate}[label=\textup{(\alph*)}]
\item an exact integer lower certificate at \(a_0\), already supplied in \cref{app:cert};
\item an exact prefix/quotient maximum-fiber certificate at \(a_1\);
\item an exact primitive split-pencil census certificate at \(a_1\), using the base-field model of \cref{prob:BC-expanded};
\item an exact primitive shift-pair certificate at \(a_1\);
\item a machine-checkable summation showing the total numerator is at most \(B_*\).
\end{enumerate}
Completing this checklist proves \cref{conj:finite-frontier} by \cref{thm:finite-adjacent-compiler}.
\end{problem}


\section{Entropy calculus for the frontier}\label{sec:entropy-calculus}

The unsafe construction is elementary, but its optimized location is most transparent in entropy coordinates.  This section records the calculation so that the finite integers in \cref{conj:finite-frontier} can be compared with the asymptotic formula.

Let \(m=(\rho+g)n\), \(K=\rho n+O(1)\), and \(w=m-K=gn+O(1)\).  The identity-prefix floor has logarithm
\[
        \log_2\binom nm-w\log_2|\BB|+O(\log n)
        = n\bigl(H_2(\rho+g)-\beta g\bigr)+O(\log n),
\]
where \(\beta=\log_2|\BB|\).  The crossing occurs where
\[
        H_2(\rho+g)=\beta g.
\]
This is \cref{eq:gstar}.  For \(g\) below the crossing, the prefix floor is exponentially large; for \(g\) above it, the average fiber is exponentially small.  The conjecture says that no primitive fiber or split-pencil cell beats this transition after quotient structure is charged to its own scale.

\begin{lemma}[uniqueness and monotonicity of the crossing]\label[lemma]{lem:gstar-unique}
Fix \(0<\rho<1\) and \(\beta>0\).  On any interval where \(\rho+g>1/2\), the function
\[
        F(g)=H_2(\rho+g)-\beta g
\]
is strictly decreasing.  Hence the frontier equation has at most one solution in the deployed rate-\(1/2\) band \(g>0\).  Moreover moving the agreement by \(R\) points beyond the crossing changes the entropy balance by
\[
        -\left(\beta-\log_2\frac{1-\rho-g}{\rho+g}\right)R+O(R^2/n).
\]
\end{lemma}

\begin{proof}
The derivative of binary entropy is
\[
        H_2'(x)=\log_2\frac{1-x}{x}.
\]
Thus
\[
        F'(g)=\log_2\frac{1-\rho-g}{\rho+g}-\beta.
\]
In the deployed band \(\rho=1/2\) and \(g>0\), the logarithm is negative and \(\beta>0\), so \(F'(g)<0\).  The expansion for an \(R\)-point movement follows from Taylor's formula with \(g\) changed by \(R/n\) and then multiplied by \(n\).
\end{proof}

\begin{corollary}[logarithmic reserve absorbs polynomial losses]\label[corollary]{cor:log-reserve}
Suppose Q, BC, and SP have asymptotic losses bounded by \(n^C\).  If the safe agreement is moved \(R\ge A\log n\) points beyond the entropy crossing, with \(A\) depending only on \(C,\rho,\beta\), then the entropy deficit beats all polynomial losses.
\end{corollary}

\begin{proof}
By \cref{lem:gstar-unique}, each additional agreement point beyond the crossing costs a positive constant number of bits in the deployed dense band.  Thus \(R=A\log n\) gives an entropy deficit \(\Omega(A\log n)\), while the polynomial loss costs at most \(C\log n\) bits.  Taking \(A\) large enough absorbs the loss.
\end{proof}

\subsection{Finite orientation}

For finite certificates, the entropy approximation is only a guide; all final comparisons are integer comparisons.  Still, the entropy orientation explains why the adjacent values in \cref{conj:finite-frontier} are so close to the asymptotic frontier.  In the MCA rows, the unsafe comparison is
\[
        \binom nm>p^{m-k-1}\left\lfloor\frac{q}{2^t}\right\rfloor,
\]
where \(t=128\) for KoalaBear and \(t=100\) for Mersenne-31.  The next-step comparison replaces \(m\) by \(m+1\) and multiplies the denominator by one more factor of \(p\), while the numerator changes by
\[
        \frac{\binom n{m+1}}{\binom nm}=\frac{n-m}{m+1}.
\]
Thus one agreement step changes the orientation by approximately
\[
        \log_2\frac{n-m}{m+1}-\log_2p,
\]
about \(-31\) bits at the deployed rows.  This explains the one-step interleaving between list and MCA rows: at the same \(m\), MCA pays prefix depth \(m-k-1\), one base-field factor less than the list row, but the next agreement step costs slightly more than one base-field factor.

\begin{proposition}[one-step list/MCA interleaving]\label[proposition]{prop:list-mca-interleave}
In the identity-prefix regime \(m>n/2\), compare the list inequality at depth \(m-k\) with the MCA inequality at depth \(m-k-1\), with the same base field and extension field.  If the list row is unsafe at agreement \(m\), then the MCA row is unsafe at agreement \(m\).  If the MCA row is unsafe at agreement \(m+1\), then the list row is unsafe at agreement \(m\) only if the extra ratio
\[
        p\frac{n-m}{m+1}
\]
still beats the target-budget difference.  In the deployed rows this leaves the list unsafe agreement exactly one below the MCA unsafe agreement.
\end{proposition}

\begin{proof}
At equal \(m\), the MCA lower floor has denominator \(p^{m-k-1}\) while the list floor has denominator \(p^{m-k}\), so the MCA comparison is one factor of \(p\) easier.  Moving from \(m\) to \(m+1\) multiplies \(\binom nm\) by \((n-m)/(m+1)\) and the denominator by \(p\).  The exact deployed conclusion is the integer comparison verified in \cref{thm:exact-unsafe}; the displayed factor gives the symbolic reason for the pattern.
\end{proof}

\subsection{Endpoint conventions revisited}

A recurring source of off-by-one errors is the difference between real radii and integer agreements.  The finite theorem uses closed Hamming balls.  If the last unsafe agreement is \(a_0\) and the first safe agreement is \(a_1=a_0+1\), then the threshold supremum is
\[
        \delta^*=\frac{n-a_0}{n}.
\]
The endpoint is unsafe because \(\ceil{(1-\delta^*)n}=a_0\).  Every \(\delta<\delta^*\) sufficiently close to the endpoint has integer agreement at least \(a_1\) and is safe.  This is the convention used in the values
\[
        \frac{981105}{2097152},
        \qquad
        \frac{981129}{2097152}
\]
for the two MCA rows.


\section{Circle and map-smooth rows}\label{sec:circle}

The prefix floors need only complete polynomial fibers, not a multiplicative group.  This is why the same mechanism applies to circle line rounds.

For \(p\equiv3\pmod4\), let \(\mathcal U\) be the circle group \(x^2+y^2=1\), cyclic of order \(p+1\), and let \(\chi(x,y)=x\).  A twin coset is \(gH\cup g^{-1}H\subseteq\mathcal U\), with \(g^2\notin H\).  Its \(x\)-coordinate image \(D\) has size \(|H|\).  The Chebyshev map \(T_c\) satisfies
\[
        T_c(\chi(u))=\chi(u^c),
\]
so dyadic subgroups give complete uniform Chebyshev fibers on \(D\).  Therefore \cref{prop:map-floor} applies verbatim with \(\varphi=T_c\).  Over a field containing \(i\), the bivariate circle code is diagonally equivalent to a Reed--Solomon code on the torus twin coset; over fields not containing \(i\), the same line-round analysis is obtained by the stereographic model and the base-field component decomposition.

\begin{proposition}[circle transfer]\label[proposition]{prop:circle-transfer}
For the Mersenne-31 line-round row with \(p'=2^{31}-1\), \(n=2^{21}\), and \(k=2^{20}\), the identity-prefix and flexible-budget certificates of \cref{thm:exact-unsafe} apply exactly as stated.  At nontrivial dyadic scales, the same graded prefix floors apply with \(X^c\) replaced by \(T_c\).
\end{proposition}

\begin{proof}
The row has a twin-coset \(x\)-domain \(D\subseteq\F_{p'}\) of size \(2^{21}\).  Identity-scale locators use only \(D\subseteq\BB\), so \cref{lem:identity-floor,prop:exact-witness-list,cor:flexible} apply without change.  At scale \(c\), the map \(T_c\) has complete uniform fibers on the dyadic twin-coset domain, so \cref{prop:map-floor} applies.  The code equivalence between the circle line round and the corresponding Reed--Solomon row preserves Hamming agreements and linear-combination slopes, hence preserves CA and MCA numerators.
\end{proof}

\section{What remains}\label{sec:remaining}

The paper leaves no ambiguity about which parts are theorem-level and which are conjectural.

\begin{itemize}[leftmargin=2em]
\item The unsafe side is unconditional: identity-prefix floors and flexible-budget conversion prove the deployed unsafe agreements in \cref{thm:exact-unsafe}; the universal field-size cap is \cref{thm:universal-cap}.
\item The elementary safe side is unconditional: deep-regime MCA is \cref{thm:deep-mca}, and the MCA-from-CA pincer is \cref{thm:mca-from-ca}.  The half-distance safe edge uses the imported CA theorem of \cite{BCIKS20}.
\item The exact Proximity Prize frontier is conditional on Q, BC, and SP.  These inputs are not cosmetic: the Mersenne-31 adjacent margins are only about three bits, so the finite theorem needs exact constants.
\end{itemize}

The most attackable current form of Q is the moment-kernel formulation of \cref{prop:moment-kernel} and the collision-gap hierarchy of \cref{prop:moment-sandwich}.  The most attackable current form of BC is the base-field split-pencil census: after removing quotient, tangent, common-GCD, and bounded chart strata, the remaining primitive determinantal incidence variety should have only model-size point count up to subexponential loss.  SP is the primitive residue left by the exact second-moment stratification of prefix collisions.

A proof of all three inputs with polynomial losses would settle the asymptotic Proximity Prize frontier with logarithmic reserve.  A proof with the finite constants fitting inside \cref{tab:cert-margins} would settle the adjacent pairs of \cref{conj:finite-frontier} exactly.

\appendix

\section{Exact arithmetic certificate format}\label{app:cert}

The exact deployed inequalities can be checked in a few lines of integer arithmetic.  The following comparisons are the certificate payload.

For KoalaBear MCA:
\[
        \binom{2^{21}}{1116047}>
        (2^{31}-2^{24}+1)^{67470}
        \left\lfloor\frac{(2^{31}-2^{24}+1)^6}{2^{128}}\right\rfloor,
\]
while the same inequality with \(1116048\) and \(67471\) is false, and
\[
        \binom{\floor{(2^{31}-2^{24}+1)^6/2^{128}}+1}{2}\,2^{20}
        <(2^{31}-2^{24}+1)^6-2^{21}.
\]
For KoalaBear list:
\[
        \binom{2^{21}}{1116046}2^{128}>
        (2^{31}-2^{24}+1)^{67470}(2^{31}-2^{24}+1)^6,
\]
with failure after replacing \((1116046,67470)\) by \((1116047,67471)\).

For Mersenne-31 MCA:
\[
        \binom{2^{21}}{1116023}>
        (2^{31}-1)^{67446}
        \left\lfloor\frac{(2^{31}-1)^4}{2^{100}}\right\rfloor,
\]
with failure at \((1116024,67447)\), and
\[
        \binom{\floor{(2^{31}-1)^4/2^{100}}+1}{2}\,2^{20}
        <(2^{31}-1)^4-2^{21}.
\]
For the Mersenne-31 list row:
\[
        \binom{2^{21}}{1116022}2^{100}>
        (2^{31}-1)^{67446}(2^{31}-1)^4,
\]
with failure at \((1116023,67447)\).

\section{Minimal verification script}\label{app:script}

The script below verifies the integer comparisons without relying on a slow
large-binomial routine.  It builds each binomial coefficient from its prime
factorization and then performs ordinary exact integer comparisons.

{\small
\begin{verbatim}
from math import isqrt

n = 2**21
k = 2**20
p_kb = 2**31 - 2**24 + 1
p_m31 = 2**31 - 1


def primes_upto(N: int) -> list[int]:
    sieve = bytearray(b"\x01") * (N + 1)
    sieve[:2] = b"\x00\x00"
    for i in range(2, isqrt(N) + 1):
        if sieve[i]:
            start = i * i
            sieve[start:N + 1:i] = b"\x00" * (((N - start) // i) + 1)
    return [i for i in range(N + 1) if sieve[i]]


def v_factorial(N: int, p: int) -> int:
    total = 0
    while N:
        N //= p
        total += N
    return total


def balanced_product(items: list[int]) -> int:
    if not items:
        return 1
    while len(items) > 1:
        nxt = []
        it = iter(items)
        for a in it:
            b = next(it, None)
            nxt.append(a if b is None else a * b)
        items = nxt
    return items[0]


def comb_exact(N: int, K: int, primes: list[int]) -> int:
    K = min(K, N - K)
    factors = []
    for p in primes:
        e = v_factorial(N, p)
        e -= v_factorial(K, p) + v_factorial(N - K, p)
        if e:
            factors.append(pow(p, e))
    return balanced_product(factors)


PRIMES = primes_upto(n)
COMB_CACHE: dict[int, int] = {}


def C(m: int) -> int:
    if m not in COMB_CACHE:
        COMB_CACHE[m] = comb_exact(n, m, PRIMES)
    return COMB_CACHE[m]


def check_mca(name: str, p: int, ext: int, bits: int, m: int) -> None:
    q = p**ext
    w = m - (k + 1)
    B = q // (2**bits)
    assert C(m) > p**w * B, f"{name} MCA pass failed"
    assert not (C(m + 1) > p**(w + 1) * B), f"{name} MCA next failed"
    assert ((B + 1) * B // 2) * k < q - n, f"{name} MCA pole failed"
    print(f"{name} MCA ok: m={m}, w={w}, target=2^-{bits}")


def check_list(name: str, p: int, ext: int, bits: int, m: int) -> None:
    q = p**ext
    w = m - k
    lhs = C(m) * 2**bits
    rhs = p**w * q
    assert lhs > rhs, f"{name} list pass failed"
    next_lhs = C(m + 1) * 2**bits
    next_rhs = p**(w + 1) * q
    assert not (next_lhs > next_rhs), f"{name} list next failed"
    print(f"{name} list ok: m={m}, w={w}, target=2^-{bits}")


if __name__ == "__main__":
    check_mca("KoalaBear", p_kb, 6, 128, 1116047)
    check_list("KoalaBear", p_kb, 6, 128, 1116046)
    check_mca("Mersenne-31", p_m31, 4, 100, 1116023)
    check_list("Mersenne-31", p_m31, 4, 100, 1116022)
\end{verbatim}
}



\section{Expanded lower-certificate proofs}\label{app:expanded-lower}

This appendix gives a slower version of the lower-certificate argument.  The body of the paper states the construction in its shortest form; here we spell out the three verifications that a referee or implementation should check: the locator-prefix list identity, the injectivity of the produced codewords, and the reason the identity scale is the asymptotic terminal obstruction.

\subsection{The prefix word and its exact list}\label{appsub:prefix-word}

Fix integers \(K\le m\le n\), and put \(w=m-K\).  For a subset \(M\subseteq D\), \(|M|=m\), write
\[
        \Lambda_M(X)=X^m+c_1(M)X^{m-1}+\cdots+c_m(M),
        \qquad c_j(M)=(-1)^j e_j(M).
\]
The prefix map is
\[
        \Phi_w(M)=(c_1(M),\ldots,c_w(M))\in\BB^w.
\]
For any prefix \(z=(z_1,\ldots,z_w)\), define
\[
        P_z(X)=X^m+z_1X^{m-1}+\cdots+z_wX^{m-w},
        \qquad U_z=P_z|_D .
\]
If \(M\in\Phi_w^{-1}(z)\), then
\[
        U_z-\Lambda_M|_D
        =-(c_{w+1}(M)X^{m-w-1}+\cdots+c_m(M))|_D .
\]
The right-hand side has degree at most \(m-w-1=K-1\).  Hence it is a codeword of \(\RS[\F,D,K]\).  It agrees with \(U_z\) at every point of \(M\), because \(\Lambda_M\) vanishes there.

The converse is equally important.  Suppose \(c=Q|_D\), \(\deg Q<K\), agrees with \(U_z\) on at least \(m\) positions.  Then
\[
        R(X)=P_z(X)-Q(X)
\]
is monic of degree \(m\), has the prescribed first \(w\) coefficients, and has at least \(m\) roots in \(D\).  Therefore its roots in \(D\) are exactly some set \(M\) of size \(m\), and \(R=\Lambda_M\).  Consequently \(M\in\Phi_w^{-1}(z)\), and \(c=U_z-\Lambda_M|_D\).  Thus the list at radius \(1-m/n\) is exactly the prefix fiber.  No hidden codewords appear.

This exactness is what makes the lower side robust.  The proof never assumes random behavior of elementary symmetric functions.  It only uses the average fiber size to guarantee that some prefix is heavy.

\subsection{Injectivity and field of definition}\label{appsub:injectivity-field}

If two subsets \(M,M'\in\Phi_w^{-1}(z)\) gave the same codeword, then
\[
        \Lambda_M|_D=\Lambda_{M'}|_D .
\]
The difference \(\Lambda_M-\Lambda_{M'}\) has degree at most \(m-1<n\), since the leading monic terms cancel.  It vanishes on all of \(D\), a set of size \(n\), hence it is the zero polynomial.  Therefore \(\Lambda_M=\Lambda_{M'}\) and \(M=M'\).  This proves that the list size is the fiber size, not merely the number of fiber elements counted with multiplicity.

When \(D\subseteq\BB\), all coefficients \(c_j(M)\) lie in \(\BB\).  The heaviest prefix can therefore be chosen in \(\BB^w\), and the received word \(U_z\) is \(\BB\)-valued.  The codewords produced from the locator tails are also \(\BB\)-valued.  In an extension row \(\BB\subseteq\F\), this base-field valued list is then converted into extension-valued slopes by choosing a pole in \(\F\setminus\BB\).  The list construction and the pole conversion use the two fields in different ways: the list mass is controlled by \(|\BB|^w\), whereas the challenge budget is controlled by \(|\F|\).

\subsection{Quotient scales and why they must be paid separately}\label{appsub:quotient-paid}

Let \(\varphi:D\to Q\) be a complete uniform scale of degree \(c\).  A quotient support is a pullback \(\varphi^{-1}(A)\), with \(A\subseteq Q\).  Its locator is
\[
        \prod_{a\in A}(\varphi(X)-a).
\]
This is a monic polynomial of degree \(c|A|\), but its coefficients are constrained by the polynomial \(\varphi\).  If \(\varphi(X)=X^c\), only powers congruent to the same residue class appear in the leading part.  If \(\varphi=T_c\) is a Chebyshev map on a circle row, the same sparsity appears in the Chebyshev basis.

These quotient floors are real lower obstructions.  A safe-side theorem that tries to estimate all supports by an aperiodic random model will fail, because quotient supports have fewer effective prefix coordinates.  The correct upper proof first descends to the quotient domain, applies the quotient version of Q there, and charges the resulting contribution to the quotient ledger.  Only supports primitive at every scale should be sent to the split-pencil census BC.

\subsection{The entropy crossing from the exact inequality}\label{appsub:entropy-crossing}

The deployed exact inequality has the form
\[
        \binom nm > |\BB|^{m-K} B_*,
        \qquad B_*\approx \epsstar |\F| .
\]
Taking base-two logarithms gives
\[
        \log_2\binom nm-(m-K)\log_2|\BB|>
        \log_2 B_* .
\]
For \(m=(\rho+g)n\), \(K=\rho n+O(1)\), and fixed extension degree, the right side is \(O(1)\) or at most \(O(\log |\BB|)\), while the left side is
\[
        n\left(H_2(\rho+g)-g\log_2|\BB|\right)+O(\log n).
\]
Thus the asymptotic crossing is determined by
\[
        H_2(\rho+g)=g\log_2|\BB|.
\]
The finite row retains the lower-order constants, the target exponent, and the extension degree.  That is why the paper states both an asymptotic frontier and exact adjacent-pair predictions.

\section{Expanded pole-conversion and extension-line checks}\label{app:pole-extension-expanded}

The pole conversion is the bridge from list decoding in \(C^+\) to MCA for \(C\).  This appendix records the exact verification and the two common pitfalls.

\subsection{Why a pole produces degree \texorpdfstring{\(<k\)}{<k} words}\label{appsub:pole-degree}

Let \(P\in\F[X]\) have degree \(<k+1\), and let \(\alpha\notin D\).  The polynomial
\[
        P(X)-P(\alpha)
\]
vanishes at \(X=\alpha\), so it is divisible by \(X-\alpha\).  The quotient
\[
        Q_\alpha(X)=\frac{P(X)-P(\alpha)}{X-\alpha}
\]
has degree \(<k\).  Therefore, if \(P\) agrees with a received word \(U\) on a support \(S\), then the line
\[
        f_\alpha(x)=\frac{U(x)}{x-\alpha},
        \qquad g_\alpha(x)=-\frac1{x-\alpha}
\]
has the property that
\[
        f_\alpha(x)+P(\alpha)g_\alpha(x)=Q_\alpha(x)
\]
on \(S\).  Each distinct value \(P(\alpha)\) is a slope explained on the support of the corresponding listed polynomial.

The far condition is equally elementary.  If \(g_\alpha\) agreed with a degree \(<k\) polynomial \(G\) on more than \(k\) points of \(D\), then
\[
        (X-\alpha)G(X)+1
\]
would be a nonzero polynomial of degree at most \(k\) with more than \(k\) roots.  It is nonzero because its value at \(\alpha\) is \(1\).  Hence \(g_\alpha\) is farther than radius \(1-(k+1)/n\), and in particular no support of size \(>k\) can jointly explain the pair by degree \(<k\) words.  This is the mutuality check.

\subsection{Collision counting and bucket numerators}\label{appsub:pole-buckets}

Let \(P_1,\ldots,P_L\) be the listed degree \(<k+1\) polynomials.  For a pole \(\alpha\), the number of bad slopes obtained is the number of distinct values among \(P_i(\alpha)\).  If a collision occurs for a pair \(i\ne j\), then \(\alpha\) is a root of \(P_i-P_j\), a nonzero polynomial of degree at most \(k\).  Hence each pair forbids at most \(k\) poles.  If
\[
        k\binom L2<|\Omega|,
\]
where \(\Omega\) is the allowed pole set, some pole has no collisions.  This is the collision-free certificate used in the deployed rows.

When the inequality is not available, one can still use buckets.  Let \(b_\alpha(y)=|\{i:P_i(\alpha)=y\}|\).  The number of equal-value pairs at \(\alpha\) is \(\sum_y\binom{b_\alpha(y)}2\).  Averaging over \(\Omega\) gives an upper bound on the average collision count, because each polynomial difference contributes at most \(k\) roots.  Cauchy--Schwarz then gives a pole with at least
\[
        \frac{L|\\Omega|}{|
        \Omega|+kL}
\]
distinct buckets.  In the body we write this as \(L(q-n)/(q-n+kL)\) when \(\Omega=\F\setminus D\).  The bucket version is useful for asymptotic profile lower bounds; the collision-free version is stronger and cleaner for the finite deployed certificates.

\subsection{Extension poles versus base-valued lines}\label{appsub:extension-pitfall}

A base-valued line \((f_1,f_2)\) with \(f_i:D\to\BB\) cannot create a genuinely new extension-slope MCA failure merely by sampling \(\gamma\in\F\setminus\BB\).  Restriction of scalars decomposes the equation
\[
        f_1+
        \gamma f_2=c
\]
into base-field component equations.  Since \(1\) and \(\gamma\) are linearly independent over \(\BB\), a support explanation of the combination recovers support explanations of both components.  This is the confinement statement in \cref{cor:confinement-core}.

The simple-pole construction avoids this inertness by making the generator itself extension-valued: \((x-\alpha)^{-1}\) with \(\alpha\notin\BB\).  The line is not projectively base-valued, and the slopes \(P_i(\alpha)\) naturally live in \(\F\).  Therefore the extension-field lower certificate is not an artifact of using a larger denominator.  It is an actual extension-valued adversarial line.

\section{Formal certificate grammar}\label{app:certificate-grammar}

This appendix gives the finite proof grammar in a form that can be translated directly into a checker.  It is intentionally more formal than the main text.  The goal is to ensure that a future proof of Q, BC, and SP has a precise landing site.

\subsection{Objects}

A row object is a tuple
\[
        \mathcal R=(n,k,\BB,\F,D,\varepsilon^*,\mathrm{type}),
\]
where \(D\subseteq\BB\), \(|D|=n\), \(\F/\BB\) is the slope field, and \(\mathrm{type}\in\{\mathrm{MCA},\mathrm{list}\}\).  A verdict object at agreement \(a\) contains:
\[
        \mathsf{Verdict}(a)=
        (\mathsf{unsafe\_lower},\mathsf{safe\_upper},\mathsf{cells},\mathsf{dependencies}).
\]
The lower field is an integer lower bound for the numerator at agreement \(a\).  The upper field is an integer upper bound.  A verdict is unsafe if \(\mathsf{unsafe\_lower}>B_*\), safe if \(\mathsf{safe\_upper}\le B_*\), and undecided otherwise.

\begin{definition}[cell record]\label[definition]{def:cell-record}
A cell record is a tuple
\[
        (\mathrm{name},\mathrm{profile},\mathrm{scope},\mathrm{bound},\mathrm{proof\_tag}).
\]
The profile contains the agreement \(a\), the code dimension \(K\), the prefix depth \(w=a-K\), the complement degree \(\omega=n-a\), and, for split-pencil cells, the shifted degrees \((d_1,d_2)\).  The scope is a predicate specifying which witnesses are assigned to the cell.  The bound is an integer.  The proof tag is one of
\[
        \mathrm{PF},\mathrm{SPole},\mathrm{Deep},\mathrm{MC},\mathrm{HJ},\mathrm{ImportCA},
        \mathrm{Q},\mathrm{BC},\mathrm{SP},\mathrm{Quot},\mathrm{Tangent},\mathrm{GCD}.
\]
\end{definition}

The proof tags mean: prefix floor, simple-pole conversion, deep-regime MCA, MCA-from-CA, half-Johnson CA, imported CA proximity gap, prefix/quotient flatness, split-pencil census, shift-pair control, quotient recursion, tangent-slope count, and common-divisor removal.

\subsection{Disjointness convention}

The finite sum must not count a witness in two primitive cells.  The checker should implement the following priority order:
\[
\mathrm{Common}\prec
\mathrm{Tangent}\prec
\mathrm{Quotient}\prec
\mathrm{Prefix}\prec
\mathrm{GCD}\prec
\mathrm{ShiftPair}\prec
\mathrm{SplitPencilPrimitive}.
\]
A witness is assigned to the first applicable class.  This convention is conservative: moving a witness earlier in the order can only decrease the burden on primitive cells.  The primitive split-pencil census BC is therefore asked only about witnesses that survive all earlier deletion predicates.

\begin{proposition}[soundness of the grammar]\label[proposition]{prop:grammar-soundness}
Assume each cell record has a valid proof tag and a true integer upper bound for its scope.  If the scopes cover all witnesses at agreement \(a\) under the priority convention, then the sum of the cell bounds is a valid numerator upper bound at agreement \(a\).
\end{proposition}

\begin{proof}
The priority convention turns the scopes into disjoint assigned scopes.  Coverage says every witness belongs to one assigned scope.  Each assigned scope has size at most its cell bound.  Summing the cell bounds therefore bounds the total number of witnesses.  In MCA rows, the witness numerator is the number of bad slopes; the assignment may further split a slope by a chosen canonical support, which only overcounts slopes by witnesses and is therefore safe for an upper bound.  In list rows, the witness numerator is a list size or a binomial moment projected to list elements; the cell compiler must state the projection factor, and the same disjointness argument applies.
\end{proof}

\subsection{Unsafe records}

An unsafe prefix-floor record has the form
\[
        \mathsf{PF}(n,K,m,p,L_0):
        \quad
        \binom nm>p^{m-K}L_0.
\]
For a list row, \(L_0=B_*\) or \(B_*+1\), depending on whether strictness is encoded before or after flooring.  For an MCA row, one uses \(K=k+1\) and \(L_0=B_*\), then attaches a simple-pole admissibility record
\[
        \mathsf{SPole}(L_0):
        \quad
        \binom{L_0+1}{2}k<q-n
\]
or the stronger non-base-pole version with \(q-|\BB|\).  The output is \(L_0+1\) bad slopes and hence probability \((L_0+1)/q>\varepsilon^*\).

\subsection{Safe records}

A safe record at agreement \(a\) has the schematic form
\[
\begin{array}{rcl}
\mathsf{safe\_upper}(a)&=&
U_{\rm tangent}(a)+U_{\rm deep}(a)+U_{\rm quotient}(a)\\
&&{}+U_Q(a)+U_{BC}(a)+U_{SP}(a).
\end{array}
\]
The first two summands are explicit from \cref{thm:deep-mca,thm:mca-from-ca}.  The quotient term is a finite recursion over complete scale maps.  The Q term is a list of maximum prefix-fiber bounds.  The BC term is a list of primitive split-pencil profile bounds.  The SP term is a primitive shift-pair bound.  The record is accepted if the sum is at most \(B_*\).

\begin{remark}[machine-checkable but not machine-proved]
The grammar does not require a computer to prove Q, BC, or SP.  It only requires that once a mathematical proof supplies their finite constants, the final assembly contains no informal judgement.  Every row then reduces to exact integer arithmetic and a finite list of tagged dependencies.
\end{remark}

\section{Expanded dependency audit}\label{app:dependency-audit}

This appendix restates the proof dependencies in a linear order.  It is designed to help a reader decide where a disagreement with the frontier prediction would have to enter.

\begin{enumerate}[leftmargin=2em,label=\textbf{D\arabic*.}]
\item \textbf{Reed--Solomon interpolation.}  The only algebraic facts used at the base level are uniqueness of a degree \(<k\) polynomial from \(k\) values, the Lagrange interpolation formula, and the fact that two degree \(<k\) polynomials agree on at most \(k-1\) points unless equal.
\item \textbf{Prefix floor.}  The map from \(m\)-subsets to leading locator prefixes has codomain \(\BB^w\).  A heaviest fiber gives a list floor.  Exact witness classification follows because a polynomial of degree \(m\) with \(m\) roots is its locator.
\item \textbf{Simple pole.}  A list in \(\RS[\F,D,k+1]\) gives slopes through evaluation at \(\alpha\notin D\).  Pairwise collision of slope values is controlled by roots of degree \(\le k\) difference polynomials.
\item \textbf{Unsafe certificates.}  The four deployed lower bounds are only the integer inequalities in \cref{app:cert}, plus simple-pole admissibility for MCA rows.
\item \textbf{Deep MCA.}  Below one third of the distance, three close explanations force all close slopes to ride one tangent pair; the number of such slopes is at most \(r+1\).
\item \textbf{MCA-from-CA.}  Below half the distance, a column-close pair contributes only tangent slopes, while a column-far pair turns every MCA witness into a CA witness.
\item \textbf{CA input.}  The paper has an elementary half-Johnson CA theorem and one imported unique-decoding CA theorem.  Removing the import weakens only the half-distance safe anchor.
\item \textbf{Slope elimination.}  At agreement \(m\), support explanation is equivalent to \(w=m-K\) interpolation functionals vanishing.  For a line, this is a rank-one condition on two \(w\)-vectors.
\item \textbf{Lattice dichotomy.}  The rank-one support census can be written as a denominator-divisibility problem in a rank-two polynomial module.  Low shifted degree is a near-rational/tangent cell; balanced shifted degree is the split-pencil residue.
\item \textbf{Q.}  Boundary split-pencil profiles are exactly prefix fibers.  The most concrete form is the maximum-fiber bound for \(\Phi_w\), or equivalently the high-moment/collision hierarchy after quotient recursion.
\item \textbf{BC.}  Interior balanced profiles are counted by a base-field-normalized split-pencil census.  The base-field normalization is forced by explicit subfield floors.
\item \textbf{SP.}  Primitive shift-pair collisions in the moment ledger are controlled separately after quotient-pullback shift pairs are deleted.
\item \textbf{Finite staircase.}  Once one unsafe agreement and the next safe agreement are certified, monotonicity determines the closed-ball threshold exactly.
\end{enumerate}

\begin{proposition}[localization of possible counterexamples]\label[proposition]{prop:counterexample-localization}
Suppose an adjacent-pair prediction in \cref{conj:finite-frontier} fails, while the exact unsafe certificate at \(a_0\) is correct.  Then at agreement \(a_0+1\), at least one of the finite Q, BC, or SP certificates must be false or incomplete.
\end{proposition}

\begin{proof}
The unsafe certificate at \(a_0\) is independent.  If the adjacent prediction fails, then \(a_0+1\) is not safe, so the true numerator exceeds \(B_*\).  The paid cells are bounded by theorem-level results listed in D1--D9.  If Q, BC, and SP finite certificates were also true and covered their scopes, \cref{prop:grammar-soundness,thm:finite-adjacent-compiler} would give a numerator at most \(B_*\), contradiction.  Thus the excess witness must either be a prefix/quotient excess, a primitive split-pencil excess, or a primitive shift-pair excess.
\end{proof}

\section{Pseudocode for the final finite scanner}\label{app:scanner-pseudocode}

The following pseudocode is not part of the mathematical proof, but it fixes the intended shape of the final checker.

{\small
\begin{verbatim}
for row in deployed_rows:
    B = floor(row.epsilon * row.q)
    a0 = row.proved_unsafe_agreement
    a1 = a0 + 1

    assert prefix_floor_passes(row, a0)
    if row.type == "MCA":
        assert simple_pole_admissible(row, a0)

    cells = []
    cells += tangent_cells(row, a1)
    cells += deep_or_CA_cells(row, a1)
    cells += quotient_cells(row, a1)
    cells += Q_certificate_cells(row, a1)
    cells += BC_certificate_cells(row, a1)
    cells += SP_certificate_cells(row, a1)

    assert covers_all_witness_profiles(cells, priority_order)
    U = sum(cell.bound for cell in cells)
    assert U <= B

    print(row.name, "first safe agreement", a1)
\end{verbatim}
}

The functions \texttt{prefix\_floor\_passes} and \texttt{simple\_pole\_admissible} are already implemented by the verification script in \cref{app:script}.  The remaining functions are the future finite Q/BC/SP certificates.  Their mathematical content is not hidden in the scanner; it is exactly the content of \cref{prob:Q,prob:BC-expanded,prob:SP-expanded}.


\section{Glossary of the three settlement inputs}\label{app:input-glossary}

This appendix gives a compact but explicit glossary for Q, BC, and SP.  It is included to prevent the common misunderstanding that the conditional theorem assumes a vague ``randomness'' statement.  Each input is a concrete counting assertion with a finite version and an asymptotic version.

\subsection*{Q: prefix/quotient flatness}

\textbf{Object.}  A complete scale map \(\varphi:D\to Q\) and the prefix map on monic locators over \(Q\),
\[
        M\longmapsto (c_1(M),\ldots,c_w(M)).
\]
At identity scale \(\varphi\) is the identity.  At quotient scale \(\varphi\) is a power map \(X\mapsto X^c\), a Chebyshev map \(T_c\), or another complete polynomial fiber map.

\textbf{Average model.}  The average fiber has size \(\binom{|Q|}{m}/|\BB|^w\).  The unsafe construction uses only the fact that the maximum fiber is at least the average.  The safe side asks that, after quotient-pullback fibers are assigned to their own scales, the maximum primitive fiber is not much larger than the average.

\textbf{Asymptotic target.}  Max-to-mean ratio \(e^{o(n)}\), or polynomial when a logarithmic reserve is used.

\textbf{Finite target.}  Exact integer maximum-fiber bounds at the four adjacent safe agreements.  In the Mersenne rows, the bound must be within about three bits of the average-scale prediction after all other cells are included.

\textbf{Equivalent forms.}  For \(w<\operatorname{char}\BB\), Q is equivalent to a constant-weight moment-curve syndrome problem.  Its collision moments count balanced \(\{0,\pm1\}\)-valued words in the moment kernel.  Boundary split-pencil profiles are exactly Q.

\subsection*{BC: base-field-normalized split-pencil census}

\textbf{Object.}  A balanced weak-Popov basis \((W_1,N_1),(W_2,N_2)\) with
\[
        W_1N_2-W_2N_1=\gamma P_D,
\]
and the divisibility count for denominators in the pencil
\[
        AW_1+BW_2\mid P_D.
\]
The profile records \(\omega=\deg(AW_1+BW_2)\), the shifted degrees \((d_1,d_2)\), and the prefix depth \(w\).

\textbf{Average model.}  There are two unavoidable scales.  The challenge-field primitive term is roughly \(\binom n\omega q^{1-w}\).  The base-field floor is \(M_{\BB}(d_1;m)\) from \eqref{eq:MB-floor}.  The correct model is the maximum of these terms and \(1\), after deleting quotient-pullback and common-GCD cells.

\textbf{Asymptotic target.}  Primitive profile counts bounded by the corrected model times \(e^{o(n)}\).

\textbf{Finite target.}  Explicit integer upper bounds for each profile at \(a_0+1\).  A profile may be certified by direct algebraic rigidity, by a finite chart atlas, by a determinant-rank certificate, or by a machine-checkable exhaustive count over a reduced finite parameter set.

\textbf{Equivalent forms.}  BC is the primitive rank-one support census after slope elimination.  In the dual-weighted syndrome calculus, it is a rank-one Hankel-window problem plus a divisor test inside a two-generator pencil.

\subsection*{SP: primitive shift-pair control}

\textbf{Object.}  Collision configurations that enter the prefix moment ledger through a shift relation but are not explained by a quotient pullback or common divisor.

\textbf{Average model.}  The random-model term is the corresponding contribution to \(\Gamma_r\) for the relevant fixed or finite moment.  Quotient-induced shift pairs are not counted in SP; they are recursively assigned to Q.

\textbf{Asymptotic target.}  Polynomial or \(e^{o(n)}\) control of the primitive shift-pair contribution.

\textbf{Finite target.}  A deletion rule plus an exact integer upper bound for the remaining primitive shift pairs at the adjacent safe agreements.

\textbf{Why separate from Q?}  Q controls maximum prefix fibers.  SP controls a residual part of the collision proof of Q.  A proof of Q by a direct maximum-fiber theorem might make SP unnecessary; a proof of Q through moments must either prove SP internally or expose it as a separate certificate.

\begin{remark}[minimal route to the asymptotic theorem]
The shortest asymptotic route is: prove fixed-moment Q for every fixed \(r\), prove primitive BC with \(e^{o(n)}\) loss, and prove that the shift-pair part of the Q proof is either quotient-recursive or primitive-subexponential.  Then \cref{prop:fixed-moments-imply-Q} gives asymptotic Q, and \cref{thm:conditional-frontier} gives the MCA frontier with logarithmic reserve.
\end{remark}

\begin{remark}[minimal route to the finite theorem]
The finite route is not the same.  It probably requires direct maximum-fiber bounds or very high moments for Q, because fixed moments do not see a three-bit adjacent margin.  It also requires BC and SP constants profile by profile.  The final arithmetic, however, is small: once those constants are supplied, \cref{app:certificate-grammar} and the script in \cref{app:script} are enough to assemble the result.
\end{remark}

\section{Finite-row upper-certificate format}\label{app:finite-upper-format}

The unsafe side already has a compact machine-checkable certificate.  The finite safe side should have the same style.  This appendix describes the expected shape of such a certificate.  It is included so that the remaining Q/BC/SP work is not an informal ``prove a good bound'' request but a concrete verifier target.

\subsection{Numerator files}\label{appsub:numerator-files}

For each row and each agreement \(a\), the verifier should maintain a numerator ledger
\[
        U(a)=U_{\rm tan}(a)+U_{\rm quo}(a)+U_Q(a)+U_{SP}(a)+U_{BC}(a).
\]
The terms have the following meanings.
\begin{description}[leftmargin=2.8em,style=nextline]
\item[\(U_{\rm tan}\)] tangent and deep-regime slopes paid by \cref{thm:deep-mca,thm:mca-from-ca};
\item[\(U_{\rm quo}\)] supports pulled back from nontrivial quotient or Chebyshev scales;
\item[\(U_Q\)] primitive identity-prefix and quotient-prefix boundary fibers;
\item[\(U_{SP}\)] primitive shift-pair second-moment residue;
\item[\(U_{BC}\)] primitive balanced split-pencil supports.
\end{description}
The finite proof at the adjacent agreement \(a_1=a_0+1\) is the integer inequality
\[
        U(a_1)\le B_*=\floor{\epsstar |\F|}.
\]
The lower certificate already proves \(L(a_0)>B_*\).  The corridor theorem then gives the exact first safe agreement.

\subsection{Q certificate}\label{appsub:Q-cert}

A finite Q certificate should list every scale \(\varphi:D\to Q\) used by the row, the active support size at that scale, and an upper bound for the largest prefix fiber.  At identity scale this is a bound of the form
\[
        \max_z |\Phi_w^{-1}(z)|\le C_Q\frac{\binom nm}{|\BB|^w},
\]
with \(C_Q\) printed as an exact rational or integer ceiling.  At quotient scale the same bound is applied to the quotient domain, with primitive pullbacks separated from composite pullbacks.  The verifier must be able to check that every support counted by Q is either primitive at the stated scale or has already been charged to a coarser quotient cell.

For the asymptotic theorem, \(C_Q=n^{O(1)}\) is acceptable with reserve.  For the finite theorem, \(C_Q\) must be small enough after all union bounds that the final numerator remains below \(B_*\).  This is the key gap between a proof that is morally correct and a proof that is prize-exact.

\subsection{BC certificate}\label{appsub:BC-cert}

A finite BC certificate should enumerate balanced profiles.  For each profile it should provide:
\begin{enumerate}[label=\textup{(\roman*)},leftmargin=2.1em]
\item the degree windows for \(W_1,W_2,A,B\);
\item the number of admissible normalized window pairs \((W_1,W_2)\);
\item a proof that quotient-pullback, common-GCD, and planted-core cases have been removed;
\item an upper bound for the number of pairs \((A,B)\) for which \(AW_1+BW_2\) is a monic divisor of \(\Lambda_D\);
\item the resulting slope numerator contribution.
\end{enumerate}
The natural model count is \(\max(1,\binom n\omega |\F|^{1-w})\), corrected by base-field floors where the line is normalized over \(\BB\).  A certificate may be analytic, combinatorial, or computational, but it must be additive over disjoint profile classes.

\subsection{SP certificate}\label{appsub:SP-cert}

SP is the seam between Q and BC.  In the second-moment expansion of prefix fibers, one encounters ordered pairs of disjoint sets with equal initial moments.  Some equalities come from quotient pullback; others come from genuine primitive shift-pair coincidences.  A finite SP certificate should identify the quotient-generated shift pairs and bound the primitive residue.  The output is an integer \(U_{SP}(a_1)\) and a proof that no shift-pair contribution is simultaneously charged to Q or BC.

This disjointness matters.  A loose proof could double-count and still be valid as an upper bound, but at the Mersenne-31 adjacent margin even small double-counting may exceed the budget.  The certificate should therefore be organized as a partition, not merely as a collection of overlapping estimates.

\section{Proof routes for Q, BC, and SP}\label{app:proof-routes}

This appendix collects proof strategies that are compatible with the structure of the paper.  They are not additional assumptions; they are possible ways to prove the named inputs.

\subsection{Fourier and moment route for Q}\label{appsub:fourier-Q}

The prefix distribution can be studied through power sums.  When \(w<\operatorname{char}\BB\), Newton identities make the first \(w\) elementary symmetric functions equivalent to the first \(w\) power sums.  For an additive character \(\psi\) and frequency vector \(\tau\in\BB^w\), the relevant Fourier coefficient is
\[
        \E_{|M|=m}\psi\left(\sum_{j=1}^w \tau_j\sum_{x\in M}x^j\right).
\]
The stable recursion in \cref{sec:Q} computes this coefficient without coefficient explosion.  A proof of Q by this route would classify large Fourier modes.  Periodic modes should correspond to quotient pullbacks; primitive modes should have enough cancellation that all fixed moments remain polynomially bounded.

The obstruction is that low moments do not directly give max-fiber bounds at finite margins.  To bridge the gap one needs either high moments, an inverse theorem saying that a large max fiber creates a detectable low or medium moment, or a compression theorem that reduces the maximum to a canonical mode whose size can be computed exactly.

\subsection{Exchange-compression route for Q}\label{appsub:compression-Q}

The prefix fibers live inside the Johnson graph on \(m\)-subsets of \(D\).  Adjacent vertices differ by exchanging one element.  The moment-kernel condition says that differences of two fiber elements are codewords in a high-distance moment code.  An exchange-compression proof would define operations that do not increase the prefix vector but move subsets toward a canonical order.  If every fiber compresses to a canonical extremal fiber without shrinking, then Q reduces to computing that extremal fiber.

Such a proof would be particularly useful in the finite adjacent rows because it targets the maximum fiber directly.  It would also explain why quotient fibers are the only large nonrandom fibers: quotient structure is stable under many exchanges, whereas primitive aperiodic structure should not be.

\subsection{Determinantal and divisor route for BC}\label{appsub:determinantal-BC}

BC asks for a count of monic divisors of \(\Lambda_D\) inside two-generator pencils.  For multiplicative rows, \(\Lambda_D=X^n-\beta\).  The divisor condition
\[
        AW_1+BW_2\mid X^n-\beta
\]
can be tested by the vanishing of a remainder, or equivalently by a resultant condition.  The primitive balanced assumption says that neither window is too small, no common factor is forced, and no quotient pullback explains the divisibility.  Under these conditions, the expected dimension is negative or zero in the frontier range, so the point count should be model-size.

Possible proof tools include determinant-rank stratification, Bezout after excluding common factors, character sums over the divisor lattice, and inverse theorems for unusually many divisors in a pencil.  The paper does not choose among these tools.  It isolates the exact statement BC must prove so that any of them can be inserted.

\subsection{Additive-combinatorial route for SP}\label{appsub:SP-route}

SP concerns pairs of small signed supports with equal low moments.  It resembles a restricted sumset or zero-sum problem on the moment curve.  Quotient pullbacks give structured solutions.  The primitive problem is to show that no other family has enough multiplicity to affect the frontier.  A plausible route is an additive-combinatorial inverse theorem: if there are too many primitive shift pairs, then the support set must correlate with a low-degree polynomial factor or quotient scale, contradicting primitiveness.

For the finite rows, the inverse theorem must be quantitative.  It must output explicit constants, not only \(o(n)\) exponents.  This is why SP is separated from Q: a second-moment proof that is asymptotically convincing may still leave too much residue for the adjacent finite margins.

\section{Endpoint conventions and common failure modes}\label{app:endpoint-failure-modes}

This appendix lists the most common ways to overclaim the frontier.

\begin{enumerate}[label=\arabic*.,leftmargin=2.1em]
\item \emph{Open versus closed balls.}  The lower certificates use closed balls.  If agreement \(a_0\) is unsafe, then radius \(1-a_0/n\) is unsafe.  Safety at \(a_0+1\) gives radii strictly below that endpoint.
\item \emph{List dimension versus MCA dimension.}  The list row uses \(K=k\).  The MCA lower row uses \(K=k+1\) and then converts to \(K=k\).  Confusing these shifts moves the agreement by one.
\item \emph{Base field versus challenge field.}  Prefix fibers are counted over \(\BB\); slopes are sampled over \(\F\).  A base-valued line is extension-inert, so extension lower certificates require extension-valued poles.
\item \emph{Quotient floors are not error terms.}  They are real lower obstructions and must be paid before applying a primitive aperiodic estimate.
\item \emph{Polynomial losses are not finite constants.}  They settle the asymptotic theorem with reserve but can exceed the finite adjacent margins.
\item \emph{Low moments are not max fibers.}  The finite Q input needs either high-moment control, an inverse theorem, or an extremal compression argument.
\end{enumerate}

These points are editorial as much as mathematical: they are the difference between a promising proof sketch and a paper whose theorem statement is checkable against the Proximity Prize rows.


\begin{thebibliography}{BCHKS25}

\bibitem[PP26]{ProximityPrize}
The Proximity Prize.
\newblock The Proximity Prize: \$1M Reed--Solomon Challenge.
\newblock Preliminary release, 2026.
\newblock \url{https://proximityprize.org/}.

\bibitem[ABF26]{ABF26}
G.~Arnon, D.~Boneh, and G.~Fenzi.
\newblock Remaining questions in list decoding and correlated agreement.
\newblock IACR Cryptology ePrint Archive, Report 2026/680, 2026.
\newblock \url{https://eprint.iacr.org/2026/680}.

\bibitem[BCHKS25]{BCHKS25}
E.~Ben-Sasson, D.~Carmon, U.~Hab{\"o}ck, S.~Kopparty, and S.~Saraf.
\newblock On proximity gaps for Reed--Solomon codes.
\newblock IACR Cryptology ePrint Archive, Report 2025/2055, 2025.
\newblock \url{https://eprint.iacr.org/2025/2055}.

\bibitem[BCIKS20]{BCIKS20}
E.~Ben-Sasson, D.~Carmon, Y.~Ishai, S.~Kopparty, and S.~Saraf.
\newblock Proximity gaps for Reed--Solomon codes.
\newblock In \emph{Proc. 61st IEEE Symposium on Foundations of Computer Science (FOCS)}, 2020.
\newblock Extended manuscript: IACR Cryptology ePrint Archive, Report 2020/654.

\bibitem[BCKL21]{BCKL21}
E.~Ben-Sasson, D.~Carmon, S.~Kopparty, and D.~Levit.
\newblock Elliptic curve fast Fourier transform (ECFFT) part I: Fast polynomial algorithms over all finite fields.
\newblock arXiv:2107.08473, 2021.

\bibitem[Cho26a]{Cho26a}
P.~Chojecki.
\newblock Capacity-edge obstructions to Reed--Solomon mutual correlated agreement over smooth multiplicative domains.
\newblock Manuscript, 2026.

\bibitem[Cho26b]{Cho26b}
P.~Chojecki.
\newblock Slack, quotient cores, and the entropy gap for smooth-domain Reed--Solomon codes: list fibers, cyclotomic rigidity, Galois-amplified collisions, and the corrected slack mutual-correlated-agreement theory.
\newblock Manuscript, merged corrected edition, 2026.

\bibitem[CS25]{CS25}
E.~Crites and A.~Stewart.
\newblock On Reed--Solomon proximity gaps conjectures.
\newblock IACR Cryptology ePrint Archive, Report 2025/2046, 2025.
\newblock \url{https://eprint.iacr.org/2025/2046}.

\bibitem[HLP24]{HLP24}
U.~Hab{\"o}ck, D.~Levit, and S.~Papini.
\newblock Circle STARKs.
\newblock IACR Cryptology ePrint Archive, Report 2024/278, 2024.
\newblock \url{https://eprint.iacr.org/2024/278}.

\end{thebibliography}

\end{document}
